% Options for packages loaded elsewhere
\PassOptionsToPackage{unicode}{hyperref}
\PassOptionsToPackage{hyphens}{url}
%
\documentclass[
]{book}
\usepackage{amsmath,amssymb}
\usepackage{lmodern}
\usepackage{iftex}
\ifPDFTeX
  \usepackage[T1]{fontenc}
  \usepackage[utf8]{inputenc}
  \usepackage{textcomp} % provide euro and other symbols
\else % if luatex or xetex
  \usepackage{unicode-math}
  \defaultfontfeatures{Scale=MatchLowercase}
  \defaultfontfeatures[\rmfamily]{Ligatures=TeX,Scale=1}
\fi
% Use upquote if available, for straight quotes in verbatim environments
\IfFileExists{upquote.sty}{\usepackage{upquote}}{}
\IfFileExists{microtype.sty}{% use microtype if available
  \usepackage[]{microtype}
  \UseMicrotypeSet[protrusion]{basicmath} % disable protrusion for tt fonts
}{}
\makeatletter
\@ifundefined{KOMAClassName}{% if non-KOMA class
  \IfFileExists{parskip.sty}{%
    \usepackage{parskip}
  }{% else
    \setlength{\parindent}{0pt}
    \setlength{\parskip}{6pt plus 2pt minus 1pt}}
}{% if KOMA class
  \KOMAoptions{parskip=half}}
\makeatother
\usepackage{xcolor}
\usepackage{color}
\usepackage{fancyvrb}
\newcommand{\VerbBar}{|}
\newcommand{\VERB}{\Verb[commandchars=\\\{\}]}
\DefineVerbatimEnvironment{Highlighting}{Verbatim}{commandchars=\\\{\}}
% Add ',fontsize=\small' for more characters per line
\usepackage{framed}
\definecolor{shadecolor}{RGB}{248,248,248}
\newenvironment{Shaded}{\begin{snugshade}}{\end{snugshade}}
\newcommand{\AlertTok}[1]{\textcolor[rgb]{0.94,0.16,0.16}{#1}}
\newcommand{\AnnotationTok}[1]{\textcolor[rgb]{0.56,0.35,0.01}{\textbf{\textit{#1}}}}
\newcommand{\AttributeTok}[1]{\textcolor[rgb]{0.77,0.63,0.00}{#1}}
\newcommand{\BaseNTok}[1]{\textcolor[rgb]{0.00,0.00,0.81}{#1}}
\newcommand{\BuiltInTok}[1]{#1}
\newcommand{\CharTok}[1]{\textcolor[rgb]{0.31,0.60,0.02}{#1}}
\newcommand{\CommentTok}[1]{\textcolor[rgb]{0.56,0.35,0.01}{\textit{#1}}}
\newcommand{\CommentVarTok}[1]{\textcolor[rgb]{0.56,0.35,0.01}{\textbf{\textit{#1}}}}
\newcommand{\ConstantTok}[1]{\textcolor[rgb]{0.00,0.00,0.00}{#1}}
\newcommand{\ControlFlowTok}[1]{\textcolor[rgb]{0.13,0.29,0.53}{\textbf{#1}}}
\newcommand{\DataTypeTok}[1]{\textcolor[rgb]{0.13,0.29,0.53}{#1}}
\newcommand{\DecValTok}[1]{\textcolor[rgb]{0.00,0.00,0.81}{#1}}
\newcommand{\DocumentationTok}[1]{\textcolor[rgb]{0.56,0.35,0.01}{\textbf{\textit{#1}}}}
\newcommand{\ErrorTok}[1]{\textcolor[rgb]{0.64,0.00,0.00}{\textbf{#1}}}
\newcommand{\ExtensionTok}[1]{#1}
\newcommand{\FloatTok}[1]{\textcolor[rgb]{0.00,0.00,0.81}{#1}}
\newcommand{\FunctionTok}[1]{\textcolor[rgb]{0.00,0.00,0.00}{#1}}
\newcommand{\ImportTok}[1]{#1}
\newcommand{\InformationTok}[1]{\textcolor[rgb]{0.56,0.35,0.01}{\textbf{\textit{#1}}}}
\newcommand{\KeywordTok}[1]{\textcolor[rgb]{0.13,0.29,0.53}{\textbf{#1}}}
\newcommand{\NormalTok}[1]{#1}
\newcommand{\OperatorTok}[1]{\textcolor[rgb]{0.81,0.36,0.00}{\textbf{#1}}}
\newcommand{\OtherTok}[1]{\textcolor[rgb]{0.56,0.35,0.01}{#1}}
\newcommand{\PreprocessorTok}[1]{\textcolor[rgb]{0.56,0.35,0.01}{\textit{#1}}}
\newcommand{\RegionMarkerTok}[1]{#1}
\newcommand{\SpecialCharTok}[1]{\textcolor[rgb]{0.00,0.00,0.00}{#1}}
\newcommand{\SpecialStringTok}[1]{\textcolor[rgb]{0.31,0.60,0.02}{#1}}
\newcommand{\StringTok}[1]{\textcolor[rgb]{0.31,0.60,0.02}{#1}}
\newcommand{\VariableTok}[1]{\textcolor[rgb]{0.00,0.00,0.00}{#1}}
\newcommand{\VerbatimStringTok}[1]{\textcolor[rgb]{0.31,0.60,0.02}{#1}}
\newcommand{\WarningTok}[1]{\textcolor[rgb]{0.56,0.35,0.01}{\textbf{\textit{#1}}}}
\usepackage{longtable,booktabs,array}
\usepackage{calc} % for calculating minipage widths
% Correct order of tables after \paragraph or \subparagraph
\usepackage{etoolbox}
\makeatletter
\patchcmd\longtable{\par}{\if@noskipsec\mbox{}\fi\par}{}{}
\makeatother
% Allow footnotes in longtable head/foot
\IfFileExists{footnotehyper.sty}{\usepackage{footnotehyper}}{\usepackage{footnote}}
\makesavenoteenv{longtable}
\usepackage{graphicx}
\makeatletter
\def\maxwidth{\ifdim\Gin@nat@width>\linewidth\linewidth\else\Gin@nat@width\fi}
\def\maxheight{\ifdim\Gin@nat@height>\textheight\textheight\else\Gin@nat@height\fi}
\makeatother
% Scale images if necessary, so that they will not overflow the page
% margins by default, and it is still possible to overwrite the defaults
% using explicit options in \includegraphics[width, height, ...]{}
\setkeys{Gin}{width=\maxwidth,height=\maxheight,keepaspectratio}
% Set default figure placement to htbp
\makeatletter
\def\fps@figure{htbp}
\makeatother
\setlength{\emergencystretch}{3em} % prevent overfull lines
\providecommand{\tightlist}{%
  \setlength{\itemsep}{0pt}\setlength{\parskip}{0pt}}
\setcounter{secnumdepth}{5}
\usepackage{booktabs}
\ifLuaTeX
  \usepackage{selnolig}  % disable illegal ligatures
\fi
\usepackage[]{natbib}
\bibliographystyle{apalike}
\IfFileExists{bookmark.sty}{\usepackage{bookmark}}{\usepackage{hyperref}}
\IfFileExists{xurl.sty}{\usepackage{xurl}}{} % add URL line breaks if available
\urlstyle{same} % disable monospaced font for URLs
\hypersetup{
  pdftitle={Modelos Matemáticos en Ecología I},
  pdfauthor={Gerardo Martín},
  hidelinks,
  pdfcreator={LaTeX via pandoc}}

\title{Modelos Matemáticos en Ecología I}
\author{Gerardo Martín}
\date{2022-08-22}

\begin{document}
\maketitle

{
\setcounter{tocdepth}{1}
\tableofcontents
}
\hypertarget{preuxe1mbulo}{%
\chapter{Preámbulo}\label{preuxe1mbulo}}

\begin{center}\includegraphics[width=20.83in]{logo} \end{center}

En el curso \textbf{Modelos matemáticos en ecología} aprenderemos a utilizar algunas herramientas matemáticas para entender procesos en ecología. Los contenidos del índice se apegan al \href{Programa-curso.pdf}{programa completo del curso}, el cual se impartirá en los \href{Horario.pdf}{horarios normales establecidos}. Para conocer cuándo, cómo y qué temas se se impartirán puedes consultar la \href{Estrategia-docente.pdf}{estrategia docente}.

\hypertarget{encuadre-de-la-materia}{%
\chapter{Encuadre de la materia}\label{encuadre-de-la-materia}}

\hypertarget{criterios-de-evaluaciuxf3n}{%
\section{Criterios de evaluación}\label{criterios-de-evaluaciuxf3n}}

Las constribuciones a cada calificación parcial serán:

\begin{itemize}
\tightlist
\item
  Asistencia (25\%)
\item
  Trabajos de clase cumplidos (50\%)
\item
  Examen (25\%)
\item
  Participación (2 puntos extra máximo)
\end{itemize}

Cabe señalar, que la asistencia no corresponderá con su presencia en las sesiones sincrónicas, sino con el cumplimiento de los trabajos de clase. La participación se medirá tanto por participación directa en las sesiones sincrónicas como por el seguimiento que uds den a la clase por correo electrónico.

\hypertarget{cuxf3mo-se-daruxe1n-las-clases}{%
\section{¿Cómo se darán las clases?}\label{cuxf3mo-se-daruxe1n-las-clases}}

Trataré de apegarnos a los tiempos de actividades sincrónicas y asincrónicas establecidos en la \href{Estrategia-docente.pdf}{estrategia docente}, pero éstos son completamente flexibles. Una estrategia que me ha funcionado en cursos anteriores ha sido la de dedicar la última actividad sicrónica de la semana a resolver dudas, es en estas sesiones en las que hay muchas posibilidades de conseguir puntos de participación.

Todos los contenidos del curso, lecturas y presentaciones, se irán añadiendo a este sitio web conforme avanza el semestre. En el Google Classroom de la materia se irán anunciando las diferentes actividades y sesiones sincrónicas con anticipación suficiente. Igualmente, los examenes y resultados serán publicados a través de esta plataforma. A petición de uds, también se publicarán aquí los videos de las sesiones sincrónicas que tengamos, en especial para aquellos temas que sean mayor interés/dificultad/importancia. Los trabajos de práctica también se publicarán en Classroom.

Para finalizar, las clases sincrónicas estarán almacenadas en la playlist de la clase contenida en mi canal de youtube.

\hypertarget{reglas-del-saluxf3n}{%
\section{Reglas del salón}\label{reglas-del-saluxf3n}}

Estas, obviamente, son particulares del modelo en línea, por lo tanto aquí van las reglas de zoom:

\begin{enumerate}
\def\labelenumi{\arabic{enumi}.}
\tightlist
\item
  Micrófonos apagados
\item
  Cámaras prendidas, exceptuando:
  2.1. Si su velocidad de internet lo dificulta
  2.2. Si tienen datos limitados
\item
  Hacer muchas preguntas
\item
  Decirme si paso algo por alto
\end{enumerate}

\hypertarget{contacto}{%
\section{Contacto}\label{contacto}}

Para reportar fallos, resolver dudas y peticiones especiales grupales o individuales por favor enviar correo electrónico a \href{mailto:gerardo.mmc@enesmerida.unam.mx}{\nolinkurl{gerardo.mmc@enesmerida.unam.mx}}.

\hypertarget{unidad-i-introducciuxf3n-a-la-modelaciuxf3n}{%
\chapter{Unidad I: Introducción a la modelación}\label{unidad-i-introducciuxf3n-a-la-modelaciuxf3n}}

\hypertarget{introducciuxf3n-al-concepto-de-modelo-matemuxe1tico}{%
\section{Introducción al concepto de modelo matemático}\label{introducciuxf3n-al-concepto-de-modelo-matemuxe1tico}}

Para comenzar a hablar de modelos matemáticos es preciso tener una idea de por qué las matemáticas forman parte de un grado en ciencias naturales. Los autores \citet{otto2011biologist}, dan cuenta de qué tan comunes son las matemáticas en la labor del biólogo/ecólogo. Hacia el año 2006, una búsqueda sistemática del material publicado en revistas científicas contentiendo la palabra clave ``math'' sugiere que había \textasciitilde900 mil artículos, la mayoría concentrados en revistas especializadas como \emph{Bulletin of Mathematical Biology}, otros en las más prestigiosas \emph{Nature} o \emph{Science}. Una mirada más cercana a las revistas especializadas \emph{Ecology and Evolution} y \emph{American Naturalist} que el 35 y 60\% respectivametne de las publicaciones utilizaron algún modelo matemático para mostrar predicciones. En resumen, la labor de un ecólogo requiere del uso frecuente de matemáticas, si no es mediante el desarrollo de algún modelo, requiere de su entendimiento.

\hypertarget{quuxe9-es-un-modelo}{%
\subsection{¿Qué es un modelo?}\label{quuxe9-es-un-modelo}}

Imaginemos a seis personas invidentes, con la tarea de encontrar qué es el objeto que está frente a ellas usando únicamente el tacto. En la parábola de los seis hombres ciegos, el objeto es un elefante, de modo que la imágen que cada uno de ellos se forma del objeto depende enteramente de la parte del elefante que están tocando.

\begin{figure}

{\centering \includegraphics[width=16.17in]{Unidad-I/elefante} 

}

\caption{La parábola de los seis hombres ciegos inspeccionando un elefante.}\label{fig:unnamed-chunk-30}
\end{figure}

Quien toque los colmillos podrá pensar que se trata de una lanza, la trompa podría tratarse de una serpiente, la cola de una cuerda, las patas troncos de árbol y el cuerpo una pared. Es evidente que todas las hipótesis presentadas después de la inspección fueron erróneas, y que cuando estudiamos al mundo lo haremos igual, con la descripción de tan sólo una fracción de éste. Un séptimo hombre que pregunte a los otros seis qué fue lo que vieron podría formarse una imágen más completa del \emph{sistema} para proponer otra hipótesis: se trata de un objeto grande reposando sobre cuatro columnas con apéndices adelante y atrás.

Al estudiar los sistemas ecológicos, al igual que los ciegos, no sabremos que estamos frente a un elefante. Los sistemas ecológicos, al igual que los elefantes, también tienen componentes interconectados, y nosotros los ecólogos somos como los hombres ciegos, sólo podemos observar ciertas partes de los ecosistemas. En nuestro trabajo entonces, aprender a identificar y proponer hipótesis sobre cómo funcionan los sistemas ecológicos.

Las hipótesis propuestas por los seis hombres representan modelos, es decir simplificaciones del mundo que nos ayudan a entenderlo. Los modelos matemáticos son simplificaciones formales del mundo utilizando otro lenguaje, las matemáticas \citep{haefner1998modeling}.

\hypertarget{cuxf3mo-construir-un-modelo}{%
\section{Cómo construir un modelo}\label{cuxf3mo-construir-un-modelo}}

¿Cómo se construye un modelo? Es una pregunta difícil de responder, pues difícilmente existe un sólo modo de hacerlo que funcione para \href{mailto:tod@s}{\nolinkurl{tod@s}}. Aquí voy a hablar de mi propia experiencia, y cómo aprendí a utilizar las matemáticas y estadística para entender los sistemas ecológicos. Los pasos generales que sigo son:

\begin{enumerate}
\def\labelenumi{\arabic{enumi}.}
\tightlist
\item
  Identificar el sistema y sus componentes
\item
  Crear un modelo conceptual, haciendo dibujos
\item
  Identificar las herramientas matemáticas disponibles para representar el modelo conceptual
\item
  Establecer objetivos claros y alcanzables
\item
  Proponer una serie de modelos alternativos
\item
  Determinar una estrategia para seleccionar uno solo de los modelos que propuse para logar el objetivo
\end{enumerate}

\begin{figure}

{\centering \includegraphics[width=8.5in]{Unidad-I/Diagrama-dearrollo} 

}

\caption{ Diagrama de mi estrategia para desarrollar modelos.}\label{fig:unnamed-chunk-31}
\end{figure}

La estrategia general que les doy puede, por el momento, sonar un tanto abstracta, pero la iremos poniendo en práctica conforme avanza el curso para que al final, esperemos, tenga más sentido que ahorita.

\hypertarget{ejemplos}{%
\subsection{Ejemplos}\label{ejemplos}}

La presente pandemia ha puesto en boca de todxs muchos aspectos de la modelación matemática y estadística. Como es de esperarse, los modelos son distintos dependiendo del objetido que se quiera conseguir con ellos:

\begin{enumerate}
\def\labelenumi{\arabic{enumi}.}
\tightlist
\item
  Determinar el número reproductivo básico (\(R_0\)) de COVID-19
\item
  Predecir la ocupación hospitalaria ante ciertas medidas de restricción de movilidad
\item
  Estimar cómo cambiaría el número reproductivo básico ante ciertas medidas de restricción
\item
  Predecir localidades geográficas donde surgirán nuevas cepas
\item
  Predecir las cualidades infecciosas o virulentas de las nuevas cepas
\end{enumerate}

Todos estos objetivos puede de cierta manera conseguirse por medio de modelos matemáticos, todos ellos con modelos muy distintos, algunos con herramientas muy diferentes unos de otros.

\hypertarget{discusiuxf3n-sobre-las-distintas-herramientas-matemuxe1ticas-empleadas-en-la-modelaciuxf3n-matemuxe1tica}{%
\section{Discusión sobre las distintas herramientas matemáticas empleadas en la modelación matemática}\label{discusiuxf3n-sobre-las-distintas-herramientas-matemuxe1ticas-empleadas-en-la-modelaciuxf3n-matemuxe1tica}}

Las matemáticas son un área de conocimiento muy extenso, por lo que hablar de todas las herramientas disponibles ¡nos podría llevar años! En este curso nos vamos a enfocar en un tipo general de modelo y de sistema ecológico: aquellos que cambian rápidamente con el tiempo, también llamados dinámicos.

Las herramientas matemáticas más comunmente utilizadas para estudiar estos sistemas son las ecuaciones diferenciales ordinarias, ecuaciones por diferencias finitas y el álgebra de matrices. Las ecuaciones diferenciales y por diferencias finitas son muy parecidas entre sí, porque nos representar y medir los cambios que sufre un sistema ecológico a través del tiempo en cada uno de sus componentes. Además es fácil representar las ideas plasmadas diagramáticamente en ellas. El álgebra de matrices, por otro lado, nos sirve para \emph{conectar} los cambios de los diferentes componentes del sistema de estudio, y así hacer cálculos sobre sus cambios \emph{como un todo}. Por supuesto que para poder utilizar las matrices de este modo, tenemos que echar mano de una serie de supuestos de \emph{linealidad}, de modo que el álgebra lineal también es una herramienta básica.

Como podemos ver, gran parte de este curso estará enfocado en \emph{entender y medir los cambios}. Te preguntarás entonces ¿cambios de qué? En ecología, estos son algunos de los cambios que nos interesa observar:

\begin{enumerate}
\def\labelenumi{\arabic{enumi}.}
\tightlist
\item
  Número de individuos de una especie
\item
  Cantidad de materia orgánica almacenada en la vegetación
\item
  Calorías consumidas por una población de aves
\item
  Probabilidad de extinción de una especie
\end{enumerate}

Cuando medimos cambios entonces, es necesario observar el estado del mundo natural en diferentes puntos a lo largo del tiempo. Por lo tanto, la modelación matemática también requiere desarrollar buenas habilidades estadísticas. Estas últimas nos permitirán entender probabilísticamente a los sistemas ecológicos a partir de los mecanismos que representamos en los modelos matemáticos. Más adelante veremos a qué nos referimos con \emph{probabilísticamente}.

Antes de entrar en estos temas más filosóficos, sin embargo, debemos que conocer algunas expresiones algebráicas y geométricas que utilizaremos rutinariamente para representar fenómenos. Algunas de las expresiones y herramientas geométricas más comunes son:

\begin{enumerate}
\def\labelenumi{\arabic{enumi}.}
\tightlist
\item
  La recta, descrita por expresiones como \(y = a + bx\)
\item
  La parábola, \(y = x^2\)
\item
  Las Cónicas, \(y = 1/x\), o \(y^2 - x^2 = 1\) y muchas otras
\item
  Las funciones trigonométricas, \(\sin \theta = \frac{O}{H}\)
\end{enumerate}

Estas expresiones pueden ser combinadas dentro de un sistema de ecuaciones. Con sistema de ecuaciones nos referimos a ecuaciones separadas que están conectadas por términos en común y que representan \emph{flujos}. Veamos un sencillo sistema de ecuaciones estructuradas por edad con crecimiento logístico:

\begin{equation}
    \dot{J} = rA - \gamma A^2 - \kappa J \label{eq:juv}
\end{equation}

\begin{equation}
    \dot{A} = \kappa J \label{eq:ad}
\end{equation}

La expresión \(rA - \gamma A^2\) es claramente una parábola, y representa el crecimiento logístico. Por otra parte \(- \kappa J\) es un término lineal que, debido a su signo negativo representa un flujo de salida del compartimento \(J\) y es una entrada al compartimento \(A\). Este modelo relativamente sencillo representa la dinámica poblacional de un bicho con estados de desarrollo juvenil \(J\) y adulto \(A\), donde el nacimiento de individuos juveniles depende el número de adultos, y el crecimiento de los adultos ocurre a medida que los juveniles maduran y entran al compartimento de adultos a una tasa \(\kappa\).

\begin{figure}

{\centering \includegraphics[width=7.29in]{Unidad-I/juve-adult} 

}

\caption{Representación gráfica del modelo de crecimiento poblacional con estados de desarrollo.}\label{fig:juve-ad}
\end{figure}

\hypertarget{relaciuxf3n-entre-las-herramientas-matemuxe1ticas-y-los-sistemas-ecoluxf3gicos}{%
\subsection{Relación entre las herramientas matemáticas y los sistemas ecológicos}\label{relaciuxf3n-entre-las-herramientas-matemuxe1ticas-y-los-sistemas-ecoluxf3gicos}}

Los sistemas ecológicos son muy variados, de modo que las herramientas matemáticas disponibles o que conocemos no siempre reflejan ciertos aspectos importantes de los sistemas ecológicos. Un ejemplo notable de ello es la variación de atributos entre individuos de una población. Supongamos que estamos modelando el crecimiento de una población (como en las ecuaciones \eqref{eq:juv} y \eqref{eq:ad}), y la tasa neta de crecimiento depende de la cantidad de adultos reproductores y del crecimiento exponencial de los gametos fértiles en cada adulto reproductor, lo cual ocurre en corales y algunas plantas.

El crecimiento poblacional como tal, puede modelarse con ecuaciones diferenciales ordinarias, mientras que representar el crecimiento de los gametos fértiles dentro de cada adulto en un sistema de ecuaciones ordinarias requiere de otras herramientas mucho más avanzadas.

\hypertarget{uso-de-los-modelos-matemuxe1ticos-en-ecologuxeda}{%
\section{Uso de los modelos matemáticos en ecología}\label{uso-de-los-modelos-matemuxe1ticos-en-ecologuxeda}}

Como ya mencionamos anteriormente, una alta proporción de las publicaciones académicas en ecología utiliza modelos matemáticos para predecir. Existen sin embargo otros usos de los modelos matemáticos que se pueden describir en términos generales como analíticos, es decir para analizar y entender cómo los diferentes componentes de un sistema afectan su coportamiento. Otras veces, los modelos se utilizan para medir las respuestas del sistema de interés ante perturbaciones, haciendo muchas simulaciones. Veamos una serie de ejemplos de estos diferentes usos de los modelos.

\hypertarget{analuxedticos-dinuxe1micas-depredador-presa}{%
\subsection{Analíticos: Dinámicas depredador-presa}\label{analuxedticos-dinuxe1micas-depredador-presa}}

En 1928, Vito Volterra \citep{volterra1928variations} desarrolló un modelo matemático para describir cómo varía el número de individuos de una especie de depredador y una presa. Como ya es conocido, este modelo replica muy bien el comportamiento oscilatorio de ambas poblaciones (Figura \ref{fig:volterra}).

\begin{figure}

{\centering \includegraphics[width=11.11in]{Unidad-I/volterra} 

}

\caption{Simulación del modelo [Lotka-Volterra](https://scipy-cookbook.readthedocs.io/items/LoktaVolterraTutorial.html).}\label{fig:volterra}
\end{figure}

El comportamiento oscilatorio resulta \emph{obvio} desde una perspectiva biológica: los depredadores no se pueden reproducir a menos que se alimenten de sus presas, pero cuando hay muchos depredadores, estos consumen las presas disminuyendo sus números, resultando en menor potencial reproductivo para el depredador. Pero matemáticamente, ¿qué hace posible estas oscilaciones? Para responder esta pregunta se necesita entender bajo qué condiciones el sistema de ecuaciones que describe el sistema es estable, es decir cuándo los cambios de ambas poblaciones son cero, es decir bajo qué situaciones el sistema como un todo no experimenta cambio alguno. Este análisis matemático tiene un resultado un tanto sorprendente, pues las condiciones bajo las que dicha condición ocurre son un número imaginario, lo que indica que el sistema como tal es inestable, y por lo tanto \textbf{siempre} tiende a oscilar. El análisis de matemático de los modelos es en su mayoría netamente simbólico y requiere de buenas bases matemáticas.

\hypertarget{simulaciones-dinuxe1mica-de-enfermedades-infecciosas}{%
\subsection{Simulaciones: Dinámica de enfermedades infecciosas}\label{simulaciones-dinuxe1mica-de-enfermedades-infecciosas}}

Los sistemas ecológicos pueden ser muy complejos, lo que resulta en sistemas de ecuaciones igualmente complejos con muchas variables de estado y parámetros, y a la vez hace que el análisis sea muy complicado o imposible. En vista de ello se puede entonces \emph{resolver} el modelo utilizando herramientas computacionales. Un ejemplo que me gusta mucho es el estudio de la leptospirosis (infección por bacteria del género \emph{Leptospira}) en ratones \emph{Mastomys} de \citet{holt2006model}. Los autores formularon un modelo relativamente simple de ecuaciones diferenciales para ver cómo las constantes de su modelo afectan el número de ratones \emph{Mastomys} y de bacterias \emph{Leptospira} en el ambiente. Ello lo consiguieron aumentando en 10\% los valores de cada una de las constantes y midiendo el cambio tras cada simulación.

\begin{figure}

{\centering \includegraphics[width=22.76in]{Unidad-I/lepto} 

}

\caption{Las barras (eje $y$) indican el cambio proporcional de las diferentes cantidades de interés (color de las barras), ante el aumento del valor de cada parámetro (eje $x$), [@holt2006model].}\label{fig:lepto}
\end{figure}

Los análisis de sensibilidad mostrados en la Figura \ref{fig:lepto} son en realidad una serie de predicciones del sistema. De hecho la predicción generalmente echa mano de modelos matemáticos, así que nos centraremos en resumir: la predicción requiere de la planeación de una serie de escenarios para identificar áreas de intervención (alteración del hábitat por ejemplo) para obtener un resultado favorable de acuerdo con el modelo.

\hypertarget{mal-uso-de-los-modelos}{%
\subsection{Mal uso de los modelos}\label{mal-uso-de-los-modelos}}

Así como existe un uso adecuado de los modelos, es muy fácil darles un mal uso, generalmente relacionado con la calidad de los datos que los respaldan, el grado de entendimiento del sistema en cuestión, y con la relación entre las herramientas matemáticas utilizadas y la naturaleza del problema. \citet{holling1978adaptive} y \citet{karplus1975place}, sumarizan los siguientes usos potenciales de los modelos en relación a la calidad de los datos (1-baja calidad, 6-alta calidad):

\begin{enumerate}
\def\labelenumi{\arabic{enumi}.}
\tightlist
\item
  Atención del público
\item
  Aprendizaje sobre el sistema
\item
  Prueba de teorías
\item
  Juegos ``qué pasaría si\ldots{}''
\item
  Recomendación de acciones
\item
  Diseño de productos
\end{enumerate}

\citet{haefner1998modeling} sumariza estos usos potenciales de los modelos de acuerdo con el nivel de entendimiento y calidad de datos en distintas áreas del conocimieto.

\begin{figure}

{\centering \includegraphics[width=17.99in]{Unidad-I/uso-modelos} 

}

\caption{Relación entre el entendimiento y calidad de los datos, y los potenciales usos de los modelos en las diferentes disciplinas científicas [@haefner1998modeling].}\label{fig:unnamed-chunk-32}
\end{figure}

\hypertarget{tipos-de-modelos}{%
\section{Tipos de modelos}\label{tipos-de-modelos}}

La clasificación de los modelos depende de la herramienta que se utiliza para desarrollarlos. En la sección ``Cómo construir un modelo'', mencioné brevemente que uno de los pasos es la formulación de un modelo ``conceptual'', en forma de un esquema o diagrama. Esta clasificación la podemos ampliar un poco:

\begin{enumerate}
\def\labelenumi{\arabic{enumi}.}
\tightlist
\item
  \textbf{Modelo conceptual o verbal}. Descripciones en lenguaje natural.
\item
  \textbf{Diagramático}. Representación gráfica (diagramas de flujo o cajas).
\item
  \textbf{Físico}. Representación a escala de sistema (sistemas hidráulicos).
\item
  \textbf{Formal}. Matemático (algebráico o sistemas de ecuaciones).
\end{enumerate}

Aquí nos centraremos en ver a mayor profundidad los modelos formales, en clasificar los modelos matemáticos. \citet{haefner1998modeling} sugiere la siguiente clave dicotómica de clasificación:

\begin{enumerate}
\def\labelenumi{\arabic{enumi}.}
\item
  \textbf{¿Las matemáticas representan de manera explícita el proceso de interés?}

  1.1. \textbf{Sí}: Modelo mecanístico

  1.2. \textbf{No}: Modelo descriptivo, fenomenológico
\item
  \textbf{¿Las matemáticas representan de manera explícita las condiciones del estado futuro del sistema?}

  2.1. \textbf{Sí}: Modelo dinámico

  2.2. \textbf{No}: Modelo estático
\item
  \textbf{¿Las matemáticas representan el tiempo de manera contínua?}

  3.1. \textbf{Sí}: Modelo conínuo

  3.2. \textbf{No}: Modelo discreto (el tiempo sólo toma valores enteros, un día, un año)
\item
  \textbf{¿Las matemáticas representan explícitamente el espacio?}

  4.1. \textbf{Sí}: Modelo espacialmente heterogéneo

  4.2. \textbf{No}: Modelo espacialmente homogéneo
\item
  \textbf{¿Las matemáticas permiten la ocurrencia de eventos aleatorios?}

  5.1. \textbf{Sí}: Modelo estocástico

  5.2. \textbf{No}: Modelo determinístico
\end{enumerate}

\hypertarget{ejemplos-1}{%
\subsection{Ejemplos}\label{ejemplos-1}}

\hypertarget{modelo-de-presa-depredador-de-lotka-volterra--volterra1928variations}{%
\subsubsection{\texorpdfstring{Modelo de presa-depredador de Lotka-Volterra \citeyearpar{volterra1928variations}}{Modelo de presa-depredador de Lotka-Volterra {[}-@volterra1928variations{]}}}\label{modelo-de-presa-depredador-de-lotka-volterra--volterra1928variations}}

\begin{figure}

{\centering \includegraphics[width=12.18in]{Unidad-I/Volterra/Volterra} 

}

\caption{Modelo diagramático basado en el modelo matemático depredador-presa de Lotka-Volterra}\label{fig:volterra-diag}
\end{figure}

Este modelo representa cómo cambia el número de liebres y linces a través del tiempo con un sistema de ecuaciones diferenciales. El cambio en los números de individuos depdende de

\begin{enumerate}
\def\labelenumi{\arabic{enumi}.}
\tightlist
\item
  Crecimiento poblacional del depredador en relación a las presas disponibles para que el depredador se alimente
\item
  Crecimiento poblacional innato de las presas
\item
  Mortalidad de presas en relación a la cantidad de depredadores
\item
  Mortalidad innata del depredador
\end{enumerate}

El modelo entonces representa el mecanismo del sistema, por lo que es un modelo \textbf{mecanístico}. Debido a que representa el cambio del número de individuos en el tiempo, es un modelo \textbf{dinámico}. Por el momento no contamos con información suficiente para decir si es contínuo o discreto, pero no hay mención del componente espacial, por lo que podemos suponer que se trata de un modelo \textbf{espacialmente homogéneo}. La solución al modelo clásico de la figura \ref{fig:volterra-diag}, muestra un patrón regular, lo cual se debe a que todos los parámetros son constantes, y por lo tanto es un modelo \textbf{determinista}.

\hypertarget{modelo-del-panal-de-abejas}{%
\subsubsection{Modelo del panal de abejas}\label{modelo-del-panal-de-abejas}}

\citet{wilenski2003honey} desarrolló un modelo para representar cómo las abejas construyen un panal, por medio de una serie de reglas que rigen el comportamiento de abejas individuales. Debido a que el mecanismo está representado en el comportamiento colectivo de las abejas, es un modelo \textbf{mecanístico}. El modelo también muestra la progresión de la construcción del panal, por lo que es un modelo \textbf{dinámico}, y debido a que cualquiera de los eventos posibles pueden ocurrir en cualquier momento es en tiempo \textbf{contínuo}. También representa la estructura del panal en unidades espaciales, por lo que es \textbf{espacialmente heterogéneo}. Finalmente, los eventos que ocurren son aleatorios, o sea que el resultado de una realización del modelo será diferente de otra, por lo que es un modelo \textbf{estocástico}. Te invito a explorar este modelo por medio de la aplicación en línea de \href{https://www.netlogoweb.org/launch\#https://www.netlogoweb.org/assets/modelslib/Sample\%20Models/Biology/Honeycomb.nlogo}{NetLogo}.

A continuación veremos a mayor profundidad los modelos deterministicos y estocásticos.

\hypertarget{modelos-determinuxedsticos-y-estocuxe1sticos}{%
\subsection{Modelos determinísticos y estocásticos}\label{modelos-determinuxedsticos-y-estocuxe1sticos}}

Imaginemos que necesitamos un librero, y para ello necesitamos medir las dimensiones del espacio donde lo colocaremos. Solamente contamos con una cinta en pulgadas, pero el carpintero, para hacer el presupuesto, necesita las medidas en centímetros. Entonces para hacer la conversión de pulgadas a centímetros utilizamos el siguiente modelo:

\begin{equation}
    cm = 2.54 \times in  \label{eq:pulgadas}
\end{equation}

donde \(in\) son las pulgadas, \(cm\) los centímetros, y \(2.54\) es la constante de conversión. Entonces ya tenemos así un modelo que siempre nos dará la misma respuesta para cada valor de entrada:

\begin{itemize}
\tightlist
\item
  \(2in \rightarrow 2 \times 2.54cm = 5.08cm\)
\item
  \(11in \rightarrow 11 \times 2.54cm = 27.94cm\)
\end{itemize}

Imaginemos ahora que la constante de conversión no fuera una sola, \(2.54\), sino cientos, o miles de posibles valores cercanos a \(2.54\), y cada vez que resolvemos el modelo \eqref{eq:pulgadas} obtenemos un valor distinto. El primero es un modelo deterministico, siempre da la misma respuesta. El segundo, es un modelo estocástico. Te preguntarás entonces ¿para qué querríamos un modelo que de respuestas inconsistentes? La principal razón es que en ecología, como en otras áreas del conocimiento, los sistemas de estudio son muy variables, y los modelos estocásticos nos sirven para representar esa variabilidad y conocer algunos de los posibles resultados. En pocas palabras, si un sistema es determinista, siempre podemos conocer su comportamiento con alta precisión, mientras que en los estocáticos, no.

\hypertarget{deterministicos}{%
\subsubsection{Deterministicos}\label{deterministicos}}

En vista de mi comentario ``\ldots{} los sistemas de estudio son muy variables\ldots{}'' tal vez te surja otra pregunta, ¿entonces por qué estudiamos los modelos deterministas? Y no hay una pregunta mmuy sencilla, pero la principal es que son más sencillos. Por otra parte, si pensamos en los modelos desde un punto de vista estadístico, se podría también sugerir que podemos utilizar los modelos determinísticos para estudiar el comportamiento ``promedio'', o para caracterizar ese componente, el determinista, de nuestro sistema de estudio. En este sentido, se considera entonces que la mayoría de los sistemas y procesos ecológicos tienen un componente determinístico y uno estocástico. La mejor manera de modelar los sistemas ecológicos, entonces, es utilizando ambos. Por lo tanto, para modelar los sistemas ecológicos con ambos métodos, necesitamos aprender a utilizarlos por separado.

Una de las grandes bondades de los modelos determinísticos es que nos permiten analizar el comportamiento del sistema de estudio. Los modelos de cambio temporal, por ejemplo, han sido muy importantes en el desarrollo de la teoría del caos, la cual sirve para explicar por qué ciertos procesos del mundo natural nunca alcanzan un nivel de estabilidad (¡dinámicas depredaror-presa!).

\hypertarget{estocuxe1sticos}{%
\subsubsection{Estocásticos}\label{estocuxe1sticos}}

A diferencia de los modelos deterministas, los estocásticos nunca dan el mismo resultado. Existe un gran variedad de herramientas computacionales y formulaciones matemáticas para conseguir este tipo de comportamiento. En la mayoría de los casos, los modelos deterministas constan de dos componentes, uno determinista, relacionado con los mecanismos propios del sistema de estudio, y uno aleatorio, encargado de generar los eventos posibles o de simular la variación en el componente determinista. Esto lo podemos ilustrar con el modelo de conversión de pulgadas a centímetros:

\[
cm = 2.54 \times in
\]
el cual es determinista. Sin embargo, podemos formular una versión estocástica, que añada una pequeña cantidad a la conversión determinística cada vez que la calculamos:

\[
cm = 2.54 \times in + \varepsilon
\]

donde \(\varepsilon\) es el \emph{error} que hemos introducido. Técnicamente para hacer esto utilizamos un generador de números aleatorios en el lenguaje de programación \textbf{R}. Veamos una serie de conversiones:

\begin{table}

\caption{\label{tab:unnamed-chunk-35}Cinco realizaciones para la conversión de dos pulgadas a centímetros con el modelo determinista y estocástico.}
\centering
\begin{tabular}[t]{c|c}
\hline
Determinista & Estocástico\\
\hline
5.08 & 5.165393\\
\hline
5.08 & 4.665356\\
\hline
5.08 & 4.590375\\
\hline
5.08 & 5.034567\\
\hline
5.08 & 5.057283\\
\hline
\end{tabular}
\end{table}

Este tipo de \emph{inconsistencias} de los modelos estocásticos tienen implicaciones muy importantes para el comportamiento de modelos como el de Lotka-Volterra, que veremos más adelante.

\hypertarget{unidad-ii-introducciuxf3n-a-los-modelos-determinuxedsticos}{%
\chapter{Unidad II: Introducción a los modelos determinísticos}\label{unidad-ii-introducciuxf3n-a-los-modelos-determinuxedsticos}}

\hypertarget{la-luxednea-recta}{%
\section{La línea recta}\label{la-luxednea-recta}}

Una línea recta se puede representar matemáticamente de varias maneras, la más común de ellas es por medio de una función. Una función, en cambio es una expresión matemática que indica una serie de operaciones (aritméticas por ejemplo), que serán aplicadas a un conjunto de números en sucesión (los números reales, p.~ej.) para producir otro conjunto. Así la, función lineal más sencilla puede ser:

\begin{equation}
    y(x) = x
\end{equation}

Lo que quiere decir que \(y\) es una función de \(x\), y cada valor de \(y\) será igual a cada valor de \(x\). En lenguaje matemático, el conjunto de números de \(x\) se llama el dominio, y \(y(x)\) el codominio. Algunas funciones lineales más complejas son:

\begin{equation}
    y(x) = ax
\end{equation}

donde \(a\) es una constante, lo que quiere decir que cada valor de \(y(x)\) será producido por el producto \(a \times x\). Finalmente, la expresión más común de una ecuación lineal es:

\begin{equation}
    y(x) = b + ax
\end{equation}

donde \(b\) también es una constante que se llama intercepto, y \(a\) es la pendiente, pues esta última determina la inclinación de la línea recta representada en el plano cartesiano.

\hypertarget{representaciones-gruxe1ficas-de-la-recta}{%
\subsubsection{Representaciones gráficas de la recta}\label{representaciones-gruxe1ficas-de-la-recta}}

Comencemos por dar un repaso del plano cartesiano. Este consiste de la representación de dos conjuntos de números reales en los que representamos una función, como una regla de correspondencia entre los valores de cada conjunto, de modo que a cada valor de \(x\) corresponde uno y sólo uno de los valores de \(y\). Cada conjunto se puede representar como una recta, y cuando éstas se disponen perperdiculares dan origen al plano.

\begin{figure}

{\centering \includegraphics{Mod-mat-ECO_files/figure-latex/unnamed-chunk-36-1} 

}

\caption{Ejemplo del plano cartesiado mostrando en rojo la correspondencia de valores para $y(x) = x$, donde los puntos corresponden a los pares de valores $(x = -1, y = -1), (-0.5, -0.5), (0, 0), (0.5, 0.5), (1, 1)$.}\label{fig:unnamed-chunk-36}
\end{figure}

Ahora veamos, con un ejemplo gráfico por qué la constante \(a\) recibe el nombre de \emph{pendiente}.

\begin{figure}

{\centering \includegraphics{Mod-mat-ECO_files/figure-latex/unnamed-chunk-37-1} 

}

\caption{Gráficas de las funciones $y(x) = 0.5 x$ (en rojo) y $y(x) = 0.5 + 0.5 x$ (azul). Las líneas punteadas en colores indican el efecto de las constantes $a$ y $b$, mientas que Δ$x$ denota el cambio de una unidad de $x$.}\label{fig:unnamed-chunk-37}
\end{figure}

Como podemos ver, de manera más general, la línea recta es generada por una regla de correspondencia muy sencilla, un conjunto de valores \(X\) son multiplicados por un escalar \(a\), a los que se les puede añadir un \emph{intercepto} \(b\). El intercepto siempre es el valor de \(y(x)\) cuando \(x = 0\).

Ya se podrán imaginar que puede existir un sinnúmero de modelos lineales distintos, por ejemplo si \(y\) es una función de un gran número de conjuntos \(X\):

\begin{equation}
y(x_1, x_2, \dots, x_{n-1}, x_n) = b + a_1 x_1 + a_2 x_2 + \dots + a_{n-1} x_{n-1} + a_n x_n
\end{equation}

Gráficamente, estamos prácticamente limitados a explorar \(y(x_1, x_2)\), pero existen muchas herramientas matemáticas para entender modelos lineales más complejos que veremos más adelante. Estas herramientas y las implicaciones generales de la simplicidad del modelo lineal lo hacen uno de los modelos más útiles con que contamos para estudiar matemáticamente fenómenos dinámicos (que cambian de estado a lo largo del tiempo), como para estudiarlos de manera probabilística (con estadística). De hecho, el modelo lineal nos puede servir para entender modelos, como la parábola y el exponencial, que geométricamente no se parecen a la línea recta.

\hypertarget{la-paruxe1bola}{%
\section{La parábola}\label{la-paruxe1bola}}

Geométricamente, la parábola es muy distinta de la recta:

\begin{figure}

{\centering \includegraphics{Mod-mat-ECO_files/figure-latex/parabola-1} 

}

\caption{Ejemplo del plano cartesiado mostrando en rojo la correspondencia de valores para $y(x) = x^2$, donde los puntos corresponden a los pares de valores $(x = -1, y = 1), (-0.5, 0.25), (0, 0), (0.5, 0.25), (1, 1)$.}\label{fig:parabola}
\end{figure}

La función matemática más sencilla que puede describir esta forma en el plano cartesiano es:

\begin{equation}
    y(x) = x^2
\end{equation}

Al igual que con los modelos lineales, las parábolas pueden representarse matemáticamente de otras formas más complejas:

\begin{itemize}
\tightlist
\item
  \(y(x) = a + bx + cx^2\)
\item
  \(y(x) = a + cx^2\)
\item
  \(y(x) = -x^2\)
\end{itemize}

Si uno es observador, se dará cuenta de que estas ecuaciones son muy similares a la recta, y es que es posible considerar que \(x^2\) se puede considerar como otro conjunto tal que \(x' = x^2\), resultando así en funciones perferctamente lineales:

\begin{itemize}
\tightlist
\item
  \(y(x) = a + bx + cx'\)
\item
  \(y(x) = a + cx'\)
\item
  \(y(x) = -x'\)
\end{itemize}

Entonce, lo que realmente distingue matemáticamente a una parábola de una recta, es la doble ocurrencia de cada valor de \(y\) para los distintos valores de \(x\):

\begin{figure}

{\centering \includegraphics{Mod-mat-ECO_files/figure-latex/parabola-2-1} 

}

\caption{Correspondencia de dos valores distintos de $x$ para el mismo valor de $y$}\label{fig:parabola-2}
\end{figure}

Cuando esto ocurre, las funciones tiene un solo punto a lo largo de todos los valores de \(x\), donde sólo va a ocurrir un valor único de \(y\). A estos puntos se les conocen como mínimos (par de valores con coordenadas \((0, 0)\) en la figura \ref{fig:parabola-2}). Los máximos, en cambio están ilustrados en la figura \ref{fig:parabola-3}, e igual corresponde al punto con coordenadas \((0, 0)\).

\begin{figure}

{\centering \includegraphics{Mod-mat-ECO_files/figure-latex/parabola-3-1} 

}

\caption{Correspondencia de valores para $y = -x^2$.}\label{fig:parabola-3}
\end{figure}

En los modelos parabólicos más sencillos (\(y(x) = x^2, y(x) = -x^2\)), los puntos mínimos y máximos siempre estarán en las coordenadas \((x = 0, y = 0)\), pero es posible alterarlos. Por ejemplo en \(y(x) = 2 + x^2\), el mínimo estará en \((x = 0, y = 2)\). Para verificarlo resuelve \(y(x) = 2 + x^2\) para \(y(0)\), sustituyendo \(x\) por \(0\).

\hypertarget{ejercicio}{%
\subsubsection{Ejercicio}\label{ejercicio}}

A estas alturas, ya te podrás haber dado cuenta de ciertas propiedades de las parábolas, contesta las siguientes preguntas:

\begin{enumerate}
\def\labelenumi{\arabic{enumi}.}
\item
  ¿Qué tipo de punto único (mínimo o máximo) habrá en las siguientes parábolas?

  1.1. \(y(x) = 2 - x^2\)
  1.2. \(y(x) = 2x - 2x^2\)
  1.3. \(y(x) = 1 - x + 3x^2\)
  1.4. \(y(x) = x/2 - 4x^2/x\)
  1.5. \(y(x) = (5x)^2 - 3x + 1\)
\end{enumerate}

\hypertarget{cuxf3nicas}{%
\section{Cónicas}\label{cuxf3nicas}}

De manera muy general, una cónica es el conjunto de soluciones para una ecuación cuadrática cuando menos en dos variables. Por lo tanto, las parábolas, vistas en la sección anterior, son un tipo de cónica. Geométricamente las cónicas surgen a partir de la intersección entre un plano y dos conos tridimensionales superpuestos en sus vértices.

\includegraphics{https://www.mathplanet.com/Oldsite/media/28029/conic.png}

Las diferentes formas geométricas generadas por las cónicas pueden ser descritas algebráicamente por las siquiente expresiones (figura \ref{fig:conicas}):

\begin{enumerate}
\def\labelenumi{\arabic{enumi}.}
\tightlist
\item
  \(x^2 + y^2 = -1\)
\item
  \(x^2 + y^2 = 0\)
\item
  \(x^2 - y^2 = 0\)
\item
  \(xy - 1 = 0\)
\item
  \(x^2 - y = 0\)
\end{enumerate}

\begin{figure}

{\centering \includegraphics[width=6.93in]{Unidad-II/conicas} 

}

\caption{Representación gráfica de las ecuaciones 1-5.}\label{fig:conicas}
\end{figure}

\hypertarget{la-curva-normal}{%
\section{La curva normal}\label{la-curva-normal}}

En el curso de probabilidad y estadística, se acordarán, se mencionó repetidamente a la distribución normal. En cuestión de matemáticas, la distribución normal puede ser representada como una ecuación:

\begin{equation}
    y(x) = \frac{1}{\sigma \sqrt{2\pi}}e^{-\frac{(x-\mu)^2}{2\sigma^2}} \label{eq:normal}
\end{equation}

donde \(\mu\) es la media aritmética, \(\sigma^2\) es la varianza y \(x\) son todos los valores de la variable que observamos. Cuando tenemos datos de un experimento y los analizamos haciendo el cálculo del promedio estimamos justamente a \(\mu\) de la ecuación \eqref{eq:normal}, y al estimar la varianza obtenemos a \(\sigma^2\) de la misma ecuación. Como es importante notar, al hacer estos cálculos estamos asumiendo que nuestros datos tienen una distribución (un histograma) que puede ser descrita por la ecuación \eqref{eq:normal}.

Como vimos en la sección de parábolas, las funciones que no son monótonas (los valores de \(y\) no aumentan o disminuyen a lo largo de \(x\)), tienen un punto máximo o mínimo. La curva normal, tiene propiedades similares, aunque debido a la presencia de la función exponencial (\(e^{\dots}\)), los valores de \(y(x)\) \textbf{siempre serán positivos}. Por lo tanto, la curva normal tiene un punto máximo que corresponde a \(y(\mu)\), o sea, \(y\) alcanza su punto máximo cuando \(x=\mu\) (Figura \ref{fig:normal}).

\begin{figure}

{\centering \includegraphics{Mod-mat-ECO_files/figure-latex/normal-1} 

}

\caption{Representación gráfica de la curva normal.}\label{fig:normal}
\end{figure}

En estadística, esta función representa en el eje de las \(y\) la \emph{densidad} de la variable \(x\). Densidad se refiere a la proporción de observaciones de \(x\) dentro de un intervalo definido de \(x\). Por lo tanto, la curva normal, representa la probabilidad de observar ese valor de \(x\), de ahí que \(\mu\) sea el valor teórico más probable. En lo que respecta a \(\sigma^2\), éste representa qué tan concentrados estarán los valores de \(x\) alrededor de \(\mu\). Altos valores de \(\sigma\) se traducirán en baja concentración alrededor de \(\mu\) (Figura \ref{fig:vari}).

\begin{figure}

{\centering \includegraphics{Mod-mat-ECO_files/figure-latex/vari-1} 

}

\caption{Efecto de  la varianza sobre la la dispersión alrededor de la media.}\label{fig:vari}
\end{figure}

\hypertarget{funciones-trigonomuxe9tricas}{%
\section{Funciones trigonométricas}\label{funciones-trigonomuxe9tricas}}

En esta sección revisaremos algunas de las funciones trigonométricas y cómo se representan en el plano cartesiano.

La trigonometría es una rama de las matemáticas que se comenzó a desarrollar a partir de la necesidad de medir distancias de manera indirecta. Por ejemplo, la distancia entre la luna y la tierra y la tierra y el sol, la distancia entre dos barcos, la distancia entre un barco y el puero más cercano. Las herramientas trigonométricas entonces se comenzaron a desarrollar utilizando la geometría de triángulos rectángulos. Aquí es preciso describir el teorema de los triángulos proporcionales, dados dos triángulos con ángulos internos idénticos y longitudes de lados diferentes, los cocientes de las longitudes de sus lados serán iguales (Figura \ref{fig:trian-prop}).

\begin{figure}

{\centering \includegraphics[width=8.97in]{Unidad-II/Trian-prop} 

}

\caption{El teorema de los ángulos proporcionales es la base para las fórmulas geométricas de las funciones trigonométricas.}\label{fig:trian-prop}
\end{figure}

De modo que sin importar las longitudes de los lados, los cocientes de las longitudes siempre serán iguales, lo cual se puede extender a todos los triángulos rectángulos, la base geométrica para la trigonometría. Aquí veremos entonces las funciones trognométricas básicas seno, coseno y tangente.

\hypertarget{nomenclatura-de-triuxe1ngulos-para-trigonometruxeda-buxe1sica}{%
\subsubsection{Nomenclatura de triángulos para trigonometría básica}\label{nomenclatura-de-triuxe1ngulos-para-trigonometruxeda-buxe1sica}}

Hay una serie de términos tradicionales que se emplean en trigonometría, los cuales utilizaremos a lo largo del curso. Se trata de los nombres que se dan a los lados y ángulos internos de un triángulo rectángulo. Por lo general, el triángulo rectángulo lo interpretaremos como una representación del plano cartesiano, de modo que los \textbf{catetos} corresponden a los ejes de las \(x\) y \(y\), y son los lados que forman un ángulo recto al cruzarse. El lado que los une en secante, es la \textbf{hipotenusa}. El ángulo que forma la hipotenusa con el eje de las \(x\) tradicionalmente recibe el nombre con la letra griega \(\theta\) (teta).

\begin{figure}

{\centering \includegraphics[width=7.86in]{Unidad-II/nomenclatura} 

}

\caption{Representación gráfica de las partes del triángulo y los nombres que reciben tradicionalmente en trigonometría plana.}\label{fig:nomen-trig}
\end{figure}

\hypertarget{las-funciones-trigonomuxe9tricas-buxe1sicas}{%
\subsubsection{Las funciones trigonométricas básicas}\label{las-funciones-trigonomuxe9tricas-buxe1sicas}}

Todas las funciones trigonométricas representan una descripción del ángulo \(\theta\) como el cociente de la longitud de dos lados del triángulo rectángulo de modo que:

\begin{enumerate}
\def\labelenumi{\arabic{enumi}.}
\tightlist
\item
  \textbf{Seno}: \(\sin(\theta) = \frac{\mathrm{Opuesto}}{\mathrm{Hipotenusa}}\)
\item
  \textbf{Coseno}: \(\cos(\theta) = \frac{\mathrm{Adyacente}}{\mathrm{Hipotenusa}}\)
\item
  \textbf{Tangente}: \(\tan(\theta) = \frac{\mathrm{Opuesto}}{\mathrm{Adyacente}}\)
\end{enumerate}

Para recordar las fórmulas de cada función podemos utilizar la mnemotecnia:

\begin{itemize}
\tightlist
\item
  \textbf{Seno}: SOH (Seno, Opuesto, Hipotenusa)
\item
  \textbf{Coseno}: CAH (Coseno, Adyacente, Hipotenusa)
\item
  \textbf{Tangente}: TOA (Tangente, Opuesto, Adyacente)
\end{itemize}

\hypertarget{representaciuxf3n-gruxe1fica-de-las-funciones-trigonomuxe9tricas}{%
\subsubsection{Representación gráfica de las funciones trigonométricas}\label{representaciuxf3n-gruxe1fica-de-las-funciones-trigonomuxe9tricas}}

Para entender por qué las funciones trigonométricas se comportan como lo hacen es útil tener en cuenta que representan el cociente de dos números que cambian uno en relación del otro. Esto queda de manifiesto en la Figura \ref{fig:catetos}. Como se puede apreciar, la hipotenusa (\(H\) en la misma figura) se mantiene constante, pues representa el radio de la circunferencia, mientras que los catetos se alargan o encojen conforme la hipotenusa se acerca a los ejes \(x\) o \(y\).

\begin{figure}

{\centering \includegraphics[width=9.53in]{Unidad-II/Func-trigo} 

}

\caption{Representación gráfica del efecto del ángulo sobre la longitud de los catetos.}\label{fig:catetos}
\end{figure}

Como es de esperarse, los valores del seno y coseno están limitados por 1 y -1 cuando la longitud del cateto es igual a la hipotenusa. La tangente, por otro lado, puede tomar valores que van de \(-\infty, +\infty\) cuando el cateto adyacente tiene longitud cero \(0\). Aquí, mediremos los ángulos en radianes, o unidades de \(\pi\), donde \(\pi = 180°\)

\begin{figure}

{\centering \includegraphics{Mod-mat-ECO_files/figure-latex/graf-trig-1} 

}

\caption{Gráfica de las funciones trigonométricas básicas.}\label{fig:graf-trig}
\end{figure}

Las funciones seno, coseno y tangente son las básicas en trigonometría. Otras funciones comunes, pasadas en la inversión de las mecionadas son la secante, cosecante y cotangente respectivamente. Al igual que otras operaciones aritméticas, para las funciones trigonométricas, existen inversiones. La inversión de una multiplicación es una divisón, de la suma es la resta, del exponente es la raíz. Las inversiones de las funciones trigonométricas reciben el nombre de Arco-\(f(\theta)\), de modo que al aplicarlas al resultado de la función obtenemos el valor del ángulo:

\begin{enumerate}
\def\labelenumi{\arabic{enumi}.}
\tightlist
\item
  \(\theta = \arcsin \left(\frac{0}{H}\right)\)
\item
  \(\theta = \arccos \left(\frac{A}{H}\right)\)
\item
  \(\theta = \arctan \left(\frac{O}{A} \right)\)
\end{enumerate}

\hypertarget{aplicaciones-de-las-funciones-trigonomuxe9tricas}{%
\subsubsection{Aplicaciones de las funciones trigonométricas}\label{aplicaciones-de-las-funciones-trigonomuxe9tricas}}

En mi experiencia, las dos funciones con más aplicaciones en modelación de sistemas ambientales son la tangente y el coseno. Para comenzar, la tangente de un ángulo tiene una relación directa con la pendiente de una recta. Si recordamos la recta está descrita por la ecuación:

\begin{equation}
    y = a + bx \label{eq:lin-2}
\end{equation}

donde \(b\) es la pendiente de la recta que se forma al graficar \(x\) y \(y\) en el plano cartesiano. En realidad \(b\) es la tangente del ánguo que se forma entre la recta y el eje de las \(x\). Por lo tanto, cuando nos encontramos ante un modelo formulado por medio de ecuaciones diferenciales, por ejemplo, lo valores que obtenemos al hacer los cálculos de ese sistema es una pendiente o tangente.

El coseno, por otro lado, nos es muy útil para forzar la oscilación de algún parámetro. Ello es necesario cuando estamos interesados en representar fenómenos cíclicos como la temperatura diaria o anual. En la práctica, podemos reemplazar la pendiente \(b\) de la ecuación \eqref{eq:lin-2} por una función coseno.

\begin{figure}

{\centering \includegraphics{Mod-mat-ECO_files/figure-latex/coseno-forz-1} 

}

\caption{Gráfica de la función $y = 2 cos(x)^2 + 0.1  x$.}\label{fig:coseno-forz}
\end{figure}

\hypertarget{la-recta-como-modelo-universal-y-sus-transformaciones}{%
\section{La recta como modelo ``universal'' y sus transformaciones}\label{la-recta-como-modelo-universal-y-sus-transformaciones}}

De los modelos vistos anteriormente, la línea recta es el más flexible y común de todos. Por su simplicidad, puede ser utilizado para representar muchos fenómenos de la naturaleza, incluso aquellos cuyo comportamiento es\ldots{} !no lineal¡

Cuando decimos no lineal, nos referimos, generalmente, a los casos en que una regla de correspondencia entre \(x\) y \(y\) representada en el plano cartesiano no forman una línea recta. Ya hemos visto que, matemáticamente, es posible concebir como lineales reglas de corresondencia que geométricamente no son lineales, de ahí que la línea recta sea tan flexible. Para que una ecuación sea considerada como lineal debe cumplir con las siguientes características:

\begin{enumerate}
\def\labelenumi{\arabic{enumi}.}
\tightlist
\item
  Sólo hay adición de variables
\item
  Sólo hay multiplicación con constantes
\item
  Las variables solo están elevadas a la potencia 1
\end{enumerate}

Con base en estas definiciones, las siguientes ecuaciones \textbf{no} son lineales

\begin{enumerate}
\def\labelenumi{\arabic{enumi}.}
\tightlist
\item
  \(xy = 1\)
\item
  \(y = a + b x^2\)
\item
  \(y = e^x\)
\end{enumerate}

Aún cuando estas expresiones no son lineales, es posible transformar algunas de ellas para que adquieran propiedades lineales. Para comenzar, el caso 1, no puede ser linealizado. Para el caso 2, ya vimos un transformación posible, si asumimos que existe una \(x'\) tal que \(x' = x^2\). Para el caso 3 hay que usar un poco más de conocimiento de dos funciones matemáticas, el logaritmo y la exponencial.

\hypertarget{el-logaritmo}{%
\subsection{El logaritmo}\label{el-logaritmo}}

Para entender cómo funciona la transformación logarítmica, revisemos qué hace el logaritmo. Supongamos que queremos encontrar el valor de la incógnita \(x\) en la siguiente ecuación:

\[10^x = 20\]
La operación necesaria para invertir el radical \(10^x\) y encontrar el valor de \(x\) es precisamente el logaritmo base 10, de modo que:

\[\log_{10} 10^x = \log_{10} 20\]
\[x \log_{10} 10 = \log_{10} 20\]
sabemos que \(\log_{10} 10 = 1\) pues \(10^1 = 10\) por lo que:

\[x = \log_{10} 20 \approx 1.3\]
Así como en esta ocasión utilizamos \(\log_{10}\) (logaritmo base 10), se puede utilizar como base cualquier número \(b > 0\), si está en el conjunto de los números reales (positivos). El logaritmo más común recibe el nombre de natural, que utiliza de base la constante \(e = 2.718282 ...\) de Euler. La convención en escritura de logaritmos es:

\begin{itemize}
\tightlist
\item
  Logaritmo natural: \(\ln\)
\item
  Logaritmo base 10: \(\log\)
\item
  Logaritmo base \(b > 0, b \in \mathbb{R}\): \(\log_b\)
\end{itemize}

Aún cuando esta es la convención más común, en la modelación, la mayoría de las veces que se utiliza \(\log\) se refiere a \(\ln\), de modo que a menos que se especifique la base, \(\log\) y \(\ln\) también significan \(\log_e\).

\hypertarget{la-funciuxf3n-exponencial}{%
\subsection{La función exponencial}\label{la-funciuxf3n-exponencial}}

Una de las funciones más comunes en aplicaciones matemáticas es la exponencial que se denota por \(\exp(x)\), que significa \(e^x\). De este modo cuando \(y\) es una función exponencial de \(x\), tenemos que:

\[y(x) = \exp(x) = e^x\]
La función exponencial puede tomar formas más complejas, por ejemplo, conteniendo una ecuación lineal:

\[y(x) = \exp(a + bx) = e^{a + bx}\]
y dadas las reglas de los exponentes, también es equivalente a:

\[y(x) = e^a e^{bx} = a' e^{bx}\]

\hypertarget{linealizando-la-funciuxf3n-exponencial}{%
\subsubsection{Linealizando la función exponencial}\label{linealizando-la-funciuxf3n-exponencial}}

Como se ha mostrado en esta sección, transformar una función exponencial en lineal es relativamente sencillo. En sentido literal, las transformaciones logaritmicas no son lineales sino log-lineales pues:

\begin{equation}
    y(x) = e^{a + bx} \label{eq:expon}
\end{equation}

se convierte en:

\begin{equation}
    \log y(x) = a + bx \label{eq:log-lin}
\end{equation}

donde el lado izquierdo es una función logarítmica y el derecho es estrictamente lineal. Matemáticamente entonces, la transformación es lineal \(\forall x \in \mathbb{R}\), puesto que \(\log y\) puede tomar cualquier valor entre \(-\infty\) y \(+\infty\). Entonces, la ecuación número 3 de la lista de ecuaciones no-lineales al principio de esta sección es:

\[\log y(x) = \ln y(x) = x\]

Esta equivalencia, como veremos más adelante y como verán en modelación estadística, es sumamente útil en la modelación de sistemas dinámicos.

\hypertarget{representaciuxf3n-gruxe1fica-de-la-funciuxf3n-exponencial}{%
\subsubsection{Representación gráfica de la función exponencial}\label{representaciuxf3n-gruxe1fica-de-la-funciuxf3n-exponencial}}

En la función exponencial, la variable \(y\) crece muy rápido. De hecho, su crecimiento en cualquier punto es proporcional al valor que tiene en ese punto. Es decir que si tomamos un punto cualquiera en \(y\) y dibujamos una tangente a ese punto, la pendiente de esa recta tendrá el mismo valor que \(y\).

\begin{figure}

{\centering \includegraphics{Mod-mat-ECO_files/figure-latex/expon-1} 

}

\caption{Representación gráfica de la función exponencial. El punto rojo representa y = exp(1.5), de modo que la pendiente de la recta tangente a la curva (en rojo) en ese punto es precisamente exp(1.5).}\label{fig:expon}
\end{figure}

Esta relación entre el valor de la curva y la pendiente tiene implicaciones importantísimas en modelación de poblaciones, pues expone las razones por las cuales el crecimiento demográfico puede ser descrito como \emph{exponencial}: el crecimiento neto de una población, el total de individuos nuevos que se incorporarán a la población, en ese instante es proporcional al número de individuos repoductivos de esa población. Es por esta razón que el número total de individuos que nacen en una población es más grande cuantos más individuos haya, a pesar de que la tasa reproductiva por unidad de tiempo sea la misma.

Veamos como ejemplo, que el crecimiento de la humanidad, desde que se pueden obtener estimaciones, es aproximadamente exponencial \ref{fig:popn-wd}.

\begin{figure}

{\centering \includegraphics[width=38.22in]{Unidad-II/popn} 

}

\caption{Población humana a lo largo del tiempo, según estimaciones de [Our world in data](https://ourworldindata.org/world-population-growth).}\label{fig:popn-wd}
\end{figure}

A diferencia del ejemplo de la figura \ref{fig:expon}, donde \(y = \exp(x)\), para representar el crecimiento de la población humana tenemos que medir la tasa reproductiva, de modo que podamos establecer la relación directa entre el tiempo y el tamaño de la población. En otras palabras, medir cuántos individuos son producidos en promedio por cada individuo presente en la población en cada unidad de tiempo:

\begin{equation}
N(t) = \exp(r \cdot t) \label{eq:poblac}
\end{equation}

donde \(r\) es la tasa reproductiva de la población y utilizamos \(N\) en lugar de \(y\) para representar al \(N\)úmero de individuos, y \(t\) en lugar de \(x\) para dejar claro que \(x\) representa tiempo. Si utilizamos el conocimiento de la subsección anterior, veremos que al aplicar el logaritmo a ambos lados de la ecuación \eqref{eq:poblac} obtenemos ¡una ecuación lineal!

\begin{equation}
    \log N = r \cdot t
\end{equation}

Es evidente que un modelo tan sencillo como este tiene serias deficiencias biológicas. Por ejemplo, cuando \(t = 0\), \(\exp(r \cdot t) = 1\), de modo que se asume que siempre habrá cuando menos un individuo presente. Esto puede o no ser cierto, por lo que una forma más generalizable incluye un intercepto para el tamaño poblacional inicial:

\begin{equation}
    N(t) = N_0 e^{r \cdot t} \rightarrow \log N(t) = \log N_0 + r \cdot t
\end{equation}

\hypertarget{representaciuxf3n-gruxe1fica-del-logaritmo}{%
\subsubsection{Representación gráfica del logaritmo}\label{representaciuxf3n-gruxe1fica-del-logaritmo}}

A diferencia de la función exponencial, donde el crecimiento de \(y\) es cada vez más \emph{rápido}, en el logaritmo es más lento (figura \ref{fig:log}).

\begin{figure}

{\centering \includegraphics{Mod-mat-ECO_files/figure-latex/log-1} 

}

\caption{Representación gráfica del logaritmo con y = log x.}\label{fig:log}
\end{figure}

\hypertarget{relevancia-de-la-recta}{%
\subsection{Relevancia de la recta}\label{relevancia-de-la-recta}}

Una vez cubiertas las funciones logarítmicas y exponenciales, hemos comenzado a vislumbrar la relevancia más amplia de la recta para fenómenos que no son lineales. En el cálculo diferencial, la recta nos da los elementos para entender el significado de la derivada, y también es el bloque fundamental para implementar métodos de integración numérica relativamente sencillos, para ecuaciones diferenciales que no se pueden integrar analíticamente.

\hypertarget{representaciuxf3n-de-las-funciones-en-3-dimensiones}{%
\section{Representación de las funciones en 3 dimensiones}\label{representaciuxf3n-de-las-funciones-en-3-dimensiones}}

Hasta ahora sólo hemos visto funciones de dos variables, y su representación en el plano cartesiano. Las funciones de tres variables, requieren ser representadas en 3d. Puesto que hay tres dimensiones, pueden ser confusas. Para comenzar, al conveción para funciones de tres variables es que hay dos dominios llamados \(x\) y \(y\), y un solo codominio llamado \(z\), por lo que \(z\) es una función \(z(x, y)\). La gráfica de estas funciones, pueden ser en efecto sólo una línea en el epacio, o un plano, o algo mucho más complejo. Para representar un plano, necesitamos una ecuación lineal como

\[z(x, y) = x + y\]
y hacer los cálculos de \(z\) para todas las combinaciones posibles de \(x\) y \(y\), por ejemplo:

\begin{table}

\caption{(\#tab:tabla.xy)Todas las combinaciones posibles de (x, y) y el valor resultante de z.}
\centering
\begin{tabular}[t]{r|r|r}
\hline
x & y & z\\
\hline
0 & 0 & 0\\
\hline
1 & 0 & 1\\
\hline
2 & 0 & 2\\
\hline
3 & 0 & 3\\
\hline
0 & 1 & 1\\
\hline
1 & 1 & 2\\
\hline
2 & 1 & 3\\
\hline
3 & 1 & 4\\
\hline
0 & 2 & 2\\
\hline
1 & 2 & 3\\
\hline
2 & 2 & 4\\
\hline
3 & 2 & 5\\
\hline
0 & 3 & 3\\
\hline
1 & 3 & 4\\
\hline
2 & 3 & 5\\
\hline
3 & 3 & 6\\
\hline
\end{tabular}
\end{table}

De

\begin{figure}

{\centering \includegraphics{Mod-mat-ECO_files/figure-latex/3d-1-1} 

}

\caption{Representación gráfica de la función $z(x, y) = x + y$.}\label{fig:3d-1}
\end{figure}

Para poder hacer esta gráfica, es necesario hacer los cáculos para todas las combinaciones posibles de \(x\) y \(y\), de modo que conocemos todos los valores posibles de \(z\) para cada par de valores del dominio. Si en la representación en dos dimensiones, necesitamos \(n\) valores, en la representación 3d necesitamos \(2 \times n^2\), lo cual produce \(n^2\) diferentes valores de \(z\).

Las funciones trigonométricas pueden representarse

\begin{figure}

{\centering \includegraphics{Mod-mat-ECO_files/figure-latex/trig-3d-1} 

}

\caption{Representación en 3d de la función trigonométrica $z(x, y) = cos(x) sin(y)$.}\label{fig:trig-3d}
\end{figure}

\hypertarget{unidad-iii-introducciuxf3n-al-cuxe1lculo-diferencial-e-integral}{%
\chapter{Unidad III: Introducción al cálculo diferencial e integral}\label{unidad-iii-introducciuxf3n-al-cuxe1lculo-diferencial-e-integral}}

\hypertarget{sucesiones}{%
\section{Sucesiones}\label{sucesiones}}

En matemáticas, de manera muy general, las sucesiones pueden ser entendidas como una función cuyo dominio \textbf{siempre} es el conjunto o un subconjunto de los números naturales, mientras que el codominio puede pertenecer a cualquier otro conjunto. En notación tradicional de funciones tenemos entonces que \(x \in \mathbb{N}\), mientras que:

\begin{equation}
y(x) \in \begin{cases}\mathbb{R}\\ \mathbb{Z}\\ \mathbb{N} \end{cases}
\end{equation}

La mayoría de las sucesiones, sin embargo que se utilizan en ecología, corresponden a:

\begin{equation}
    f: 
    \begin{aligned}[t] 
            \mathbb{N}_0 \rightarrow \mathbb{R} \\
            n \rightarrow f(n)
    \end{aligned}
\end{equation}

una función \(f\) de los números naturales, incluido el cero (\(\mathbb{N}_0\)) que produce números reales (\(\mathbb{R}\)), cuyo dominio es \(n\), y codominio es \(f(n)\).

Otra característica de las sucesiones es que sus elementos están ordenados, es decir, cada elemento en el dominio y codominio tiene un sitio particular, de modo que si los conjuntos \(a_n = \{A, B, C\}\) y \(a_m = \{C, A, B\}\), \(a_n \neq a_m\), puesto que, aunque contengan los mismos elementos \(A, B\) y \(C\), estos tienen posiciones diferentes dentro de \(a_n\) y \(a_m\).

\hypertarget{notaciuxf3n-de-sucesiones}{%
\subsection{Notación de sucesiones}\label{notaciuxf3n-de-sucesiones}}

Una sucesión de los números naturales suele nombrarse como \(\{a_n\}\), seguida de la regla de correspondencia entre el conjunto \(n\) y \(a\). Así por ejemplo:

\begin{equation} 
    \{a_n\} = 2n \label{eq:a2n}
\end{equation}

es una sucesión de los números naturales pares, tal que:

\[a_1 = 1, a_2 = 4, a_3 = 6, a_4 = 8 \dots\]
Esta notación, sin embargo no es universal, por ejemplo:

-\((a_n) = a_1, a_2, a_3, \dots\)

-\((a_k)_{k=1}^m = a_1, a_2, a_3, \dots, a_m\)

-\(\{a_n\}_{n \in \mathbb{N}} = a_1, a_2, a_3, \dots\)

\hypertarget{tipos-de-sucesiones}{%
\subsection{Tipos de sucesiones}\label{tipos-de-sucesiones}}

La sucesión anterior es un ejemplo muy sencillo y poco útil en la modelación ecológica, pero es un punto de partida para entender la notación y el concepto. Existen sucesiones, sin embargo, que tienen implicaciones importantes para el desarrollo de modelos dinámicos, como las \textbf{sucesiones por recurrencia}.

\hypertarget{sucesiones-recurrentes}{%
\subsubsection{Sucesiones recurrentes}\label{sucesiones-recurrentes}}

En estas sucesiones, la regla de correspondencia indica algún tipo de operación con elementos de \(n\) en posiciones previas. Así por ejemplo, la sucesión:

\begin{equation}
\{a_n\} = a_{n-1} + a_{n-2} \label{eq:fibonacci}
\end{equation}

se conoce como sucesión de Fibonacci, en la que, el elemento \(n\) se obtiene sumando los elementos \(n-1\) y \(n-2\) de la secuencia \(\{a_n\}\).

\begin{table}

\caption{\label{tab:fibo-tabla}Secuencia de Fibonacci para los diez primeros números naturales.}
\centering
\begin{tabular}[t]{r|r}
\hline
n & a\\
\hline
1 & 1\\
\hline
2 & 1\\
\hline
3 & 2\\
\hline
4 & 3\\
\hline
5 & 5\\
\hline
6 & 8\\
\hline
7 & 13\\
\hline
8 & 21\\
\hline
9 & 34\\
\hline
10 & 55\\
\hline
\end{tabular}
\end{table}

La secuencia de finonacci, tiene un comportamiento parecido a la función exponencial, en el sentido de que el crecimiento de \(a\) es más rápido conforme los valores de \(a\) son mayores.

\begin{figure}

{\centering \includegraphics{Mod-mat-ECO_files/figure-latex/fibo-fig-1} 

}

\caption{Representación gráfica de la secuencia de Fibonacci.}\label{fig:fibo-fig}
\end{figure}

\hypertarget{sucesiones-recurrentes-muxe1s-avanzadas}{%
\subsubsection{Sucesiones recurrentes más avanzadas}\label{sucesiones-recurrentes-muxe1s-avanzadas}}

Las sucesiones recurrentes pueden ser también concebidas como modelos dinámicos en tiempo discreto, donde el tiempo está representado por el conjunto de los números naturales, así, por ejemplo, el crecimiento de una población lo podemos representar como:

\begin{equation}
    N_{t} = r N_{t-1} \label{eq:susc-crec}
\end{equation}

Como podemos ver, en \eqref{eq:susc-crec}, el crecimiento por unidad de tiempo es lineal con respecto del número de individuos en el periodo de tiempo inmediatamente previo (\(rN_{t-1}\)).

\begin{figure}

{\centering \includegraphics{Mod-mat-ECO_files/figure-latex/susc-crec-1} 

}

\caption{Simulación de la sucesión de crecimiento poblacional en tiempo discreto, para 10 unidades de tiempo, con r = 1.5, y una población inicial de N = 1.}\label{fig:susc-crec}
\end{figure}

La sucesión descrita por la ecuación \eqref{eq:susc-crec} es fácil de simular con una computadora, con aplicaciones tan populares como las hojas de cáculo como excel.

Con este tipo de sucesiones, es posible comenzar a utilizar conceptos básicos de cálculo. Podemos por ejemplo, tomar la misma sucesión de la ecuación \eqref{eq:susc-crec}, con una población inicial de \(N = 50\), y \(r = 0.6\), y simular por un periodo de tiempo de 20, para ver hacia qué valor se aproxima \(N_t\) (\ref{fig:susc-decrec}).

\begin{figure}

{\centering \includegraphics{Mod-mat-ECO_files/figure-latex/susc-decrec-1} 

}

\caption{Simulación para un valor inicial de N = 50, por 20 unidades de tiempo, y r = 0.6}\label{fig:susc-decrec}
\end{figure}

En cada simulación, \(N_t\) cambia de acuerdo con el valor de \(N_{t-1}\). Como es evidente, si \(r > 1\), \(N_t\) crece, y mientras \(t \rightarrow \infty\), \(N_t \rightarrow \infty\). En el caso contrario, cuando \(r < 1\) y \(t \rightarrow \infty\), \(N_t \rightarrow 0\)

En estos casos, decimos que por ejemplo para \(r < 1\), \(N\) está acotada por \(0\). Hay otras sucesiones, por ejemplo la de crecimiento logístico, que está acotada por la capacidad de carga, incluso si \(r > 0\).

\begin{figure}

{\centering \includegraphics{Mod-mat-ECO_files/figure-latex/susc-logis-1} 

}

\caption{Simulación para un valor inicial de N = 1, por 50 unidades de tiempo, r = 0.5, y capacidad de carga K = 50}\label{fig:susc-logis}
\end{figure}

La sucesión para el crecimiento logístico en tiempo discreto es:

\begin{equation}
N_{t} = N_{t-1} + r N_{t-1} \left(1 - \frac{N_{t-1}}{K} \right) \label{eq:susc-logis}
\end{equation}

\hypertarget{ejercicio-1}{%
\subsubsection{Ejercicio}\label{ejercicio-1}}

Utilizando la hoja de cálculo para la sucesión de crecimiento poblacional en tiempo discreto, adaptala para simular el crecimiento poblacional probando los valores de \(r\) dese 0.5 hasta 3.5 en incrementos de 0.25, y grafícalos. ¿Qué pasa con \(N_t\) con los valores de \(r\) probados?

\hypertarget{continuidad-y-luxedmites}{%
\section{Continuidad y límites}\label{continuidad-y-luxedmites}}

La continuidad y los límites son la puerta de entrada a la derivación y por lo tanto a las ecuaciones diferenciales. Cuando decimos que una función es contínua en un intervalo definido del dominio \((x_1, x_2)\), quiere decir que \(f(x)\) está definida para todos y cada uno de los posibles valores de \(x\) en ese intervalo. Existen muchas funciones que no son contínuas, por ejemplo:

\begin{equation}
    f(x) = \frac{x - 1}{x - 1} \label{eq:disc-1}
\end{equation}

Como es evidente, dicha ecuación está definida para todos los valores posibles de \(x\), excepto \(x = 1\), porque \(f(1) = \frac{1 - 1}{1-1} = \frac{0}{0}\), lo cual es una indeterminación, no existe.

\begin{figure}

{\centering \includegraphics[width=41.67in]{Unidad-III/Func-discont} 

}

\caption{Ejemplo de función con discontinuidad.}\label{fig:discont}
\end{figure}

Hay funciones con discontinuidades, gráficamente más evidentes que \eqref{eq:disc-1}, como la hipérbola, \(f(x) = 1/x\) (figura \ref{fig:hiperb}) o \(g(x) = |x|/x\), o aquellas con condicionales.

\begin{figure}

{\centering \includegraphics[width=41.67in]{Unidad-III/Hiperbola} 

}

\caption{La hipérbola.}\label{fig:hiperb}
\end{figure}

\begin{figure}

{\centering \includegraphics[width=41.67in]{Unidad-III/Deriv-absol} 

}

\caption{Cociente de valor absoluto.}\label{fig:absol}
\end{figure}

Como resulta geométricamente evidente, todas estas funciones son discontínuas para algún valor de \(x\), sin embargo existe una diferencia muy importante entre ellas y es que el valor que pueden tomar en torno al punto en que son discontínuas depende de qué lado nos acerquemos a este punto. Para la equación \(f(x) = \frac{x-1}{x-1}\), resulta claro que si nos acercamos a \(f(1)\) desde un valor \(x>1\), llegamos a un valor similar de \(f\) que si nos acercamos a \(f(1)\) por medio de un valor de \(x < 1\). Esto no sucede con las función \(g(x) = |x|/x\) puesto que si nos acercamos a \(g(0)\) por el lado positivo de \(x\) nos acercamos a \(g = 1\), mientras que, por el lado negativo, nos acercamos a \(g = -1\).

En este momento entonces es necesario comenzar a hablar del concepto de \textbf{límite}.

\hypertarget{el-luxedmite}{%
\subsection{El límite}\label{el-luxedmite}}

En el sentido estricto de la palabra, límite es el valor al que se acerca una función para un valor específico del dominio. En matemáticas, existe una notación específica para decir esto:

\begin{equation}
\lim \limits_{x \rightarrow a} x + 1
\end{equation}

donde \(\lim\) indica que necesitamos estimar el límite de la función, cuando su dominio \(x\) toma valores cada vez más cercanos a algún valor específico \(a\) (\(x \rightarrow a\)). Aquí es necesario especificar que cuando hablamos de límites \textbf{siempre} nos acercamos al valor de interés \(a\) tanto por medio de valores de \(x>a\) como \(x<a\), es decir alrededor de un intervalo de valores de \(x\) que comunmente se denomina como \(a \pm \delta\). Entonces, el límite de \(f(x)\) existe si todos los valores de \(f(a \pm \delta)\) convergen alrededor de un mismo punto.

Revisa la aplicación de Geogebra \href{https://www.geogebra.org/classic/b76f7gpx}{Límites} para visualizar este fenómeno. Juega con los valores de \(\delta\) y \(a\) para ver cómo están definidos los valores de la función alrededor de \(a\) para la función de la línea recta. Posteriormente, cambia la ecuación por:

\begin{enumerate}
\def\labelenumi{\arabic{enumi}.}
\tightlist
\item
  \(f(x) = \frac{|x|}{x}\)
\item
  \(f(x)= \frac{1}{x}\)
\end{enumerate}

para examinar qué pasa con \(f(x)\) cuando \(x \rightarrow 0\)

En los casos en que los valores de \(f(x)\) tienden a un valor único para \(f(a + \delta)\) y \(f(a - \delta)\), el límite de la función sí existe (figura \ref{fig:discont}). En cambio, como ocurre para funciones como \(1/x\) y \(|x|/x\), existe más de un valor al que ambas funciones se aproximan (figuras \ref{fig:hiperb} y \ref{fig:absol}) dependiendo se si se evalúa \(f(a + \delta)\) ó \(f(a - \delta)\), el límite no existe.

\begin{figure}

{\centering \includegraphics[width=21.38in]{Unidad-III/Limite-x} 

}

\caption{Ejemplo gráfico de una función con límite definido para cualquier valor de $x$. El valor $a$ al que $x$ se aproxima está marcado en el eje $x$ con una cruz roja, y el tamaño de ð alrededor de $a$ está delimitado con las marcas rosas. Nota cómo todos los valores posibles de $f(x)$ alrededor de $f(a)$ convergen en el mismo punto correspondiente a $f(a)$.}\label{fig:limite-def}
\end{figure}

\begin{figure}

{\centering \includegraphics[width=19.21in]{Unidad-III/Limite-abs-x} 

}

\caption{Ejemplo gráfico de una función sin límite definido cuando x→0 (equis tiende a cero). Nota cómo hay dos valores posibles a los que la función se acerca cuando x→0, ya sea cuando uno se acerca a 0 por el lado positivo de $x$ o el lado negativo de $x$.}\label{fig:limite-indef}
\end{figure}

Aún cuando estas funciones no satisfacen la definición estricta de límite, es posible encontrar los límites de la función para \(f(a)\), dependiendo del lado por el que \(x \rightarrow a\). Veamos el ejemplo de \(f(x) = |x|/x\) dado en la figura \ref{fig:limite-indef}. El límite que nos interesa es

\[ \lim \limits_{x \rightarrow 0} \frac{|x|}{x}\]
Por lo que sabemos hay dos maneras que acercarnos a \(0\), una por el lado positivo, \(x \rightarrow 0^{+}\), y otra por el lado negativo \(x \rightarrow 0^-\). Así entonces podemos establecer dos límites:

\begin{enumerate}
\def\labelenumi{\arabic{enumi}.}
\tightlist
\item
  \(\lim \limits_{x \rightarrow 0^+} \frac{|x|}{x} = 1\)
\item
  \(\lim \limits_{x \rightarrow 0^-} \frac{|x|}{x} = -1\)
\end{enumerate}

\hypertarget{cuxf3mo-se-encuentran-los-luxedmites}{%
\subsubsection{¿Cómo se encuentran los límites?}\label{cuxf3mo-se-encuentran-los-luxedmites}}

Enncontrar límites puede ser problemático, por lo que existe una variedad de estrategias cuyo uso depende del tipo de expresión que estemos evaluando. Por lo general, si tenemos \(\lim\limits_{x \rightarrow a} f(x)\) se puede comenzar por evaluar \(f(a)\), con lo que tenemos tres posibilidades:

\begin{enumerate}
\def\labelenumi{\arabic{enumi}.}
\tightlist
\item
  \(f(a) = \frac{b}{0}, b \neq 0\) (una asíntota vertical como en \(1/x\))
\item
  \(f(a) = b\) (límite encontrado)
\item
  \(f(a) = \frac{0}{0}\) (indeterminación)
\end{enumerate}

En el caso \textbf{1}, es probable que el límite no exista, lo cual depende de que \(\lim\limits_{x \rightarrow a^-} f(x) = \lim\limits_{x \rightarrow a^+} f(x)\). Para el caso \textbf{2}, el límite está claramente definido y no hay mayor problema. En el caso \textbf{3}, es necesario encontrar una manera equivalente de representar a \(f(a)\).

\hypertarget{representaciuxf3n-equivalente-de-las-funciones}{%
\subsubsection{Representación equivalente de las funciones}\label{representaciuxf3n-equivalente-de-las-funciones}}

Las tres estrategias más comunes para resolver límites donde \(f(a) = \frac{0}{0}\) son:

\begin{enumerate}
\def\labelenumi{\arabic{enumi}.}
\tightlist
\item
  Factorizar \(f(a)\)
\item
  Conjugar \(f(a)\)
\item
  Usar identidades trigonométricas
\end{enumerate}

Cabe aclarar que muchas de las transformaciones aquí descritas surten efecto bajo el supuesto de que \(x \neq a\), aunque una vez transformadas las funciones en efecto las evaluamos asumiendo que \(x = a\). Con esto en mente veamos un ejemplo para cada estrategia.

\hypertarget{factorizaciuxf3n}{%
\paragraph{Factorización}\label{factorizaciuxf3n}}

Factorizar consiste en representar un polinomio como una serie de productos de binomios (generalmente). Así, por ejemplo:

\begin{equation}
    \lim \limits_{x \rightarrow -3} \frac{x^2 + x -6}{x + 3}
\end{equation}

por sustitución directa tenemos que

\begin{equation}
\frac{(-3)^2 - 3 - 6}{-3 + 3} = \frac{0}{0}
\end{equation}

Y una vez factorizado, asumiendo que \(x \neq 3\) tenemos

\begin{equation}
 \frac{x^2 + x - 6}{x + 3} = \frac{(x + 3)(x - 2)}{x + 3} = x-2
 \end{equation}

por lo tanto, utilizando sustitución directa en la forma factorizada tenemos que:

\begin{equation}
 \lim\limits_{x \rightarrow 3} x - 2 = - 3 - 2 = -5
\end{equation}

\hypertarget{conjugaciuxf3n-racional}{%
\paragraph{Conjugación racional}\label{conjugaciuxf3n-racional}}

Una expresión conjugada de un binomio está basada en la identidad:

\[ a^2 - b^2 =(a + b)(a - b)\]

Esta identidad se puede extender a los radicales:

\[ a - b = (\sqrt{a} + \sqrt{b})(\sqrt{a} - \sqrt{b})\]

De modod que el conjugado de \(a + b\) es \(a-b\), y el de \(\sqrt{a} + \sqrt{b}\) es \(\sqrt{a} - \sqrt{b}\).

La aplicación de conjugados entonces es particularmente útil cuando queremos evaluar un el límite de una funcion racional con radicales. Por ejemplo en:

\begin{equation}
\lim \limits_{x \rightarrow 0} \frac{\sqrt{1 + x}-1}{x}
\end{equation}

por sustitución directa obtenemos \frac{0}{0}, pero no es posible factorizar el numerador, por lo tanto podemos utilizar el conjugado de \(\sqrt{1 + x} - 1\). La manera de hacerlo es, primero identificar el cojugado como:

\begin{equation}
    \sqrt{1 + x} + 1
\end{equation}

Y multiplicar tanto el nuemerados como denominador por el conjugado:

\begin{equation}
  \frac{\sqrt{1 + x}-1}{x} \cdot  \frac{\sqrt{1 + x} + 1}{\sqrt{1 + x} + 1} = \frac{1 + x-1}{x(\sqrt{1 + x} + 1)} = \frac{1}{\sqrt{1 + x} + 1}
\end{equation}

Con lo que ya podemos evaluar:

\begin{equation}
\lim\limits_{x \rightarrow 0} \frac{1}{\sqrt{1 + x} + 1} = \frac{1}{1 + 1} = \frac{1}{2}
\end{equation}

\hypertarget{identidades-trigonomuxe9tricas}{%
\paragraph{Identidades trigonométricas}\label{identidades-trigonomuxe9tricas}}

Cuano es necesario evaluar el límite de una función trigonométrica, es útil echar mano de las identidades, además de las ya bien conocidas, las pitagóricas:

\begin{enumerate}
\def\labelenumi{\arabic{enumi}.}
\tightlist
\item
  \(\mathrm{sen}^2 \theta + \cos^2 \theta = 1\)
\item
  \(\sec^2 \theta - \tan^2 \theta = 1\)
\item
  \(\csc^2 \theta - \cot^2 \theta = 1\)
\end{enumerate}

Comencemos por ver un ejemplo basado en las identidades basadas en los cocientes.

\begin{equation}
\lim \limits_{x \rightarrow 0} \frac{3x \tan x}{\mathrm{sen}x}
\end{equation}

por sustitución directa tenemos que:

\begin{equation}
\frac{3(0)\tan 0}{\mathrm{sen} 0} = \frac{0}{0}
\end{equation}

lo que indica que necesitamos representar a la función de una manera alternativa, echando mano de las identidades. Sabemos que:

\begin{equation}
\tan \theta = \frac{\mathrm{sen} \theta}{\cos \theta}
\end{equation}

por lo tanto podemos representar la función como:

\begin{equation}
\frac{3x \tan x}{\mathrm{sen}x} = \frac{3x \mathrm{sen} x}{\cos^2 x} = \frac{3x}{\cos x}
\end{equation}

Entonces, tenemos el límite que puede ser evaluado por sustitución directa:

\begin{equation}
\lim \limits_{x \rightarrow 0} \frac{3x}{\cos x} = \frac{3(0)}{\cos 0} = \frac{0}{1} = 0
\end{equation}

\hypertarget{equivalencias-de-luxedmites}{%
\subsubsection{Equivalencias de límites}\label{equivalencias-de-luxedmites}}

Los límites, como hemos visto, son números reales, por lo que todos los axiomas que ya conocemos aplican a los números reales:

\begin{enumerate}
\def\labelenumi{\arabic{enumi}.}
\item
  \(\lim \limits_{x \rightarrow a} f(x) + g(x) = \lim \limits_{x \rightarrow a} f(x) + \lim \limits_{x \rightarrow a} g(x)\)
\item
  \(\lim \limits_{x \rightarrow a} f(x) \cdot g(x) = \lim \limits_{x \rightarrow a} f(x) \cdot \lim \limits_{x \rightarrow a} g(x)\)
\item
  \(\lim \limits_{x \rightarrow a} \frac{f(x)}{g(x)} = \frac{\lim \limits_{x \rightarrow a} f(x) }{ \lim \limits_{x \rightarrow a} g(x)}\)
\item
  \(\lim \limits_{x \rightarrow a} c f(x) = c \lim \limits_{x \rightarrow a} f(x)\)
\item
  \(\lim \limits_{x \rightarrow a} f(x)^2 = \left( \lim \limits_{x \rightarrow a} f(x) \right)^2\)
\end{enumerate}

\hypertarget{derivaciuxf3n}{%
\section{Derivación}\label{derivaciuxf3n}}

Hasta el momento hemos mencionado en repetidas ocasiones una serie de conceptos relacionados muy íntimamente con la derivación. En esencia, en cálculo, el proceso de derivación se utiliza para estimar la pendiente de una recta tangente a un punto cualquiera que forma parte de una función (Figura \ref{fig:deriv}).

\begin{figure}

{\centering \includegraphics{Mod-mat-ECO_files/figure-latex/deriv-1} 

}

\caption{Representación gráfica de la tangente a un punto en la curva exponencial.}\label{fig:deriv}
\end{figure}

El proceso de estimar la pendiente de una tangente a la función en un punto \((x, f(x))\) consiste en estimar el límite:

\begin{equation}
\lim \limits_{h \rightarrow 0} \frac{f(x + h) - f(x)}{x - (x + h)} \label{eq:derivada}
\end{equation}

Resolverlo no requiere de un conocimiento profundo para encontrar que por sustitución directa de \(h = 0\), se obtiene una undeterminación (\(0/0\)). Sin embargo para entender mejor de qué se trata el límite de la ecuación \eqref{eq:derivada} podemos explorarlo con la \href{https://www.geogebra.org/classic/rjttr8kq}{aplicación para derivadas}.

\hypertarget{ejemplos-2}{%
\subsection{Ejemplos}\label{ejemplos-2}}

\hypertarget{ecuaciuxf3n-cuadruxe1tica}{%
\subsubsection{Ecuación cuadrática:}\label{ecuaciuxf3n-cuadruxe1tica}}

\[ y = x^2\]
Con lo que la derivada sería:

\begin{equation}
\lim \limits_{h \rightarrow 0} \frac{(x + h)^2 - x^2}{(x+h) - x}
\end{equation}

Utilizando las técnicas para resolver límites podemos que con este último se puede utilizar factorización para obtener:

\begin{equation}
\lim \limits_{h \rightarrow 0} \frac{(x + h)^2 - x^2}{(x+h) - x} = \lim \limits_{h \rightarrow 0} \frac{((x + h) - x)((x + h) - x)}{(x + h) + x} = \lim \limits_{h \rightarrow 0} (x + h) + x = \lim \limits_{h \rightarrow 0} 2x + h
\end{equation}

Con este resultado podemos ahora sí sustituir \(h\) directamente para ver que la derivada \(y' = 2x\).

\hypertarget{hiperbola}{%
\subsubsection{Hiperbola}\label{hiperbola}}

\[ y = \frac{1}{x}\]
Derivada:

\begin{equation}
\lim \limits_{h \rightarrow 0} \frac{\frac{1}{x + h} - \frac{1}{x}}{(x+h) - x}
\end{equation}

Podemos comenzar por simplificar las fracciones:

\begin{equation}
\lim \limits_{h \rightarrow 0} \frac{\frac{x - (x + h)}{x(x + h)}}{h} = \lim \limits - \frac{h}{x^2 + xh} \frac{1}{h} = \lim \limits_{h \rightarrow 0} -\frac{1}{x^2}
\end{equation}

Por lo tanto \(y' = \frac{1}{x^2}\)

\hypertarget{notaciuxf3n}{%
\subsection{Notación}\label{notaciuxf3n}}

La técnica de derivación se le atribuye a dos científicos, Isaac Newton y Gottfried Leibniz. La notación que se tradicionalmente enseña para comenzar con cáculo es la de Newton:

\begin{equation}
f'(x) = \lim \limits_{h \rightarrow 0} \frac{f(x + h) - f(x)}{h} \label{eq:Newton}
\end{equation}

la cual adolece de claridad, pues a menos que uno conozca bien la fórmula de la derivada no especifica el significado de \(f'(x)\). La notación de Leibniz en cambio, sí lo especifica:

\begin{equation}
\frac{df}{dx} = f'(x) \label{eq:Leibniz}
\end{equation}

donde el numerador \(df = \Delta f = f(x + h) - f(x)\) y el denominador \(dx = \Delta x = (x+h) - x = h\). La letra \(d\) se utiliza en lugar de \(\Delta\) para indicar que se trata de un cambio contínuo, en lugar de discreto, el uso tradicional de \(\Delta\).

\hypertarget{reglas-de-derivaciuxf3n}{%
\subsection{Reglas de derivación}\label{reglas-de-derivaciuxf3n}}

\hypertarget{potencias}{%
\subsubsection{Potencias}\label{potencias}}

Como ya se habrán dado cuenta, la primera regla de la derivación está relacionada con los exponentes. Por ejemplo, la derivadas de \(f(x) = x^2\) y \(f(x) = \frac{1}{x}\), son \(\frac{df}{dx} = 2x\) y \(\frac{df}{dx} = - \frac{1}{x^2}\) respectivamente. Nota, que en ambos casos, los exponentes de cada \(x\), se convirtieron en coeficientes en \(\frac{df}{dx}\), y que restamos \(1\) a cada exponente de \(x\):

\begin{enumerate}
\def\labelenumi{\arabic{enumi}.}
\tightlist
\item
  \(x^2 \rightarrow 2x^{2-1} = 2x^1 = 2x\)
\item
  \(1/x = x^{-1} \rightarrow -1x^{-1-1} = -x^{-2} = - 1/x^2\)
\end{enumerate}

Entonces la primera regla de la derivación, se puede aplicar a muchas funciones. Si \$f(x) = x\^{}n \$, \$\frac{df}{dx} = nx\^{}\{n-1\} \$.

En casos donde \(x\) aparece más de una vez, se considera que:

\[f(x) = f_1(x) + f_2{x} + \dots + f_n(x)\]

Y dado que la derivada es un límite, sabemos que \(f'(x) = f'_1(x) + f'_2(x) + \dots + f'_n(x)\).

Por ejemplo, si tenemos la función

\[y = a + bx + cx^2 + dx^3 = ax^0 + bx^1 + cx^2 + dx^3\]

su derivada será

\[ \frac{dy}{dx} = 0ax^{0-1} + bx^{1-1} + 2cx^{2-1} + 3dx^{3-1}= b + 2cx + 3dx^2\]

\hypertarget{productos}{%
\subsubsection{Productos}\label{productos}}

Si \(h(x) = f(x)g(x)\) entonces:
\[h'(x) = f'(x)g(x) + f(x)g'(x)\]

\hypertarget{cocientes}{%
\subsubsection{Cocientes}\label{cocientes}}

Si \(h(x) = \frac{f(x)}{g(x)}\), entonces:

\[h'(x) = \frac{f'(x)g(x) - f(x)g'(x)}{g(x)^2}\]

\hypertarget{funciones-en-cadena}{%
\subsubsection{Funciones en cadena}\label{funciones-en-cadena}}

Una función en cadena es aquella en que \(h(x) = (f \circ g)(x) = f(g(x))\), entonces:

\[h'(x) = f'(g(x))g'(x) = (f \circ g)'(x) g'(x)\]

Aquí es preciso dar un ejemplo:

\[h(x) = (3x-1)^2\]

donde está claro que \(g(x) = 3x^2-1\) y que \(f(u) = u^2\). Por lo tanto si evaluamos \(f'(u)\) cuando \(u = 3x^2-1\), entonces \(f'(u) = 2(3x^2-1)\). Por el otro lado, tenemos que, aplicando la regla de las potnecias, \(g'(x) = 6x\). Entonces:

\[h'(x) = 2(3x^2-1)6x\]

\hypertarget{representaciuxf3n-geomuxe9trica-de-la-derivada}{%
\subsection{Representación geométrica de la derivada}\label{representaciuxf3n-geomuxe9trica-de-la-derivada}}

Como resultado del proceso de derivación de una función \(f(x)\) se obtiene una segunda función \(f'(x)\) ó \(\frac{d}{dx}f\) cufo valor corresponde a la pendiente de una tangente para cada par de valores \((x, f(x))\).

La aplicación \href{https://www.geogebra.org/classic/w3spjhc7}{Interpretación de la derivada}, muestra ambas funciones y la recta tangente a \((x,f(x))\). Observa la pendiente de esta tangente y las coordenadas \((x, f'(x))\).

\hypertarget{aplicaciones}{%
\subsection{Aplicaciones}\label{aplicaciones}}

Uno de los modelos más generalizables que se utilizan en ecología es el de crecimiento exponencial:

\[N(t) = e^{rt}\]

Como vimos anteriormente, este modelo se puede linealizar utilizando el logaritmo:

\begin{equation}
\ln N(t) = rt \label{eq:log-N}
\end{equation}

y si resolvemos para \(r\), obtenemos que:

\[r = \frac{\ln N}{t}\]

la pendiente, o tasa de crecimiento de una poblaciones equivalente a \(\ln N/t\), para cualquier período de tiempo. Volvamos sin embargo a la ecuación \eqref{eq:log-N}, diferenciándola. Para esto utilizaremos la diferenciación implícita, pues necesitamos conocer:

\begin{enumerate}
\def\labelenumi{\arabic{enumi}.}
\tightlist
\item
  \(\frac{d}{dt}rt = r\)
\item
  \(\frac{d}{dt}N(t) \ln N = \frac{d}{dt} \frac{1}{N}\)
\end{enumerate}

por lo tanto:

\[\frac{dN}{dt} \frac{1}{N} = t\]

Con lo que obtenemos la siguiente ecuación diferencial:

\begin{equation}
\frac{dN}{dt} = rN
\end{equation}

De este modo podemos ver que, el crecimiento instantáneo de una población durante el intervalo \(dt\) es equivalente a la tasa de crecimiento exponencial \(r\). Más adelante veremos por qué \(f(x) = \ln x\), \(f'(x) = 1/x\), por el momento esta equivalencia nos sirve para ver que si \(N(t) = e^{rt}\):

\[\frac{dN}{dt} = e^{rt} = rN\]
Y si separamos ambos lados \(N\) y \(t\), tenemos:

\[\frac{dN}{N} = r \cdot dt\]

Por lo tanto, la derivada nos sirve para estudiar de cerca la tasa de crecimiento en relación al cambio de población \(dN\) y de tiempo \(dt\), con un modelo lineal.

\hypertarget{integraciuxf3n}{%
\section{Integración}\label{integraciuxf3n}}

La integración es el proceso de estimación del área bajo la curva de una función, y para nosotros como ecólogxs es muyimportante, pues la mayoría de los modelos dinámicos se formulan como ecuaciones diferenciales, los cuales debemos integrar para obtener las soluciones.

Al igual que con límites y derivadas, hay toda una notación particular de la integración:

\begin{equation}
    \int_a^b \!f(x) dx
\end{equation}

Y se lee, la integral de \(f(x)\) entre \(a\) y \(b\). \([a, b]\) corresponde a dos valores, mínimo (\(a\)) y máximo (\(b\)), de \(x\) entre los cuales se evaluará el área. El símbolo \(\int\), es una ``S'', e indica la suma de valores contínuos. Al igual que en la derivada, el signo \(d\) es la versión contínua de \(\Delta\), en las integrales \(\int\) es la versión contínua de \(\sum\).

Otra similitud entre la dericada y la integral es que est última es un límite:

\begin{equation}
    \int_a^b \!f(x) dx = \lim \limits_{n \rightarrow \infty} \sum_{i = 1}^n f(x_i) \Delta x; \Delta x = \frac{b-a}{n} \label{eq:integ-limite}
\end{equation}

Para comprederla, enfoquémonos en el lado derecho, que comenzando con el operador de sumatoria \(\sum\) indica que para cada valor \(x_i\), se evalúa la función \(f\), y se multiplica por la distancia que hay entre \(x_i\) y \(x_{i-1}\) (\(\Delta x\)). Si pensamos esto en el plano cartesiano, \(f(x_i)\) es la distancia entre \(0\) y el valor de \(y\) en \(f(x_i)\), por lo que al multiplicarlo por \(\Delta x\) estamos obteniendo el área de un rectángulo con lados de longitud \(\Delta x\) y \(f(x_i)\). Nota que conforme \(n \rightarrow \infty\), la longiud de \(\Delta x\) disminuye tanto, que el número de rectángulos con área \(\Delta x \times f(x_i)\) tiende a infinito.

\begin{figure}

{\centering \includegraphics[width=12.4in]{Unidad-III/Esquema-integral} 

}

\caption{Representación esquemática de la integral. Nota quq utilizando los rectángulos, es posible aproximar el valor de una integral con rectángulos que pasen por arriba o abajo de la curva.}\label{fig:La-integral}
\end{figure}

Puedes explorar estos conceptos con la \href{https://www.geogebra.org/classic/wmpqmr4a}{aplicación para el área bajo una recta de Geogebra}.

\hypertarget{uxe1rea-bajo-una-recta}{%
\subsection{Área bajo una recta}\label{uxe1rea-bajo-una-recta}}

Un ejemplo clásico de integración es el área bajo una recta con pendiente \(m \neq 0\) e intercepto \(0\), la cual corresponde al área de un triángulo (figura \ref{fig:Area-recta}).

\begin{figure}

{\centering \includegraphics[width=25.43in]{Unidad-III/Area-bajo-recta} 

}

\caption{Área bajo la recta f(x) = x, entre el punto 0 y a = 1.}\label{fig:Area-recta}
\end{figure}

Utilizando el método de los rectángulos tenemos que el área sombreada en la figura \ref{fig:Area-recta} entre los puntos \(0\) y \(a\) se puede aproximar sumando rectángulos con altura \(f(x)\) y ancho \(\Delta x\).

\begin{figure}

{\centering \includegraphics[width=16.04in]{Unidad-III/Rectangulos} 

}

\caption{Método de los rectángulos para calcular área bajo curva}\label{fig:Recta-area-rect}
\end{figure}

Para automatizar este proceso y encontrar todos los valores de \(x\) en el intervalo \(0, a\), podemos comenzar por tomar en cuenta el número de rectángulos, como en la figura \ref{fig:Recta-area-rect}, cuantos más rectángulos tenemos el error de cálculo será menor.

La siguiente consideración será si tomaremos en cuenta el punto inferior o superior del intervalo (figura \ref{fig:La-integral}). Si tomamos en cuenta el punto inferior del intervalo, la altura del primer rectángulo será \(f(0)\), y el ancho será \(\Delta x = a/n\). Por otra parte, tenemos que los valores de \(x\) en que evaluaremos \(f(x)\) corresponden a los valores de \(x: 0, a/n, 2a/n, 3a/n, \dots, (n-1)a/n, a\), Si comenzamos con \(n = 10\), y \(a = 1\) tenemos que evaluar \(f(x)\) en los valores:

\[x = \left \{ 0, \frac{1}{10}, \frac{1}{10}, \dots, \frac{9}{10}, 1 \right \}\]

Estas integrales se comoces como de Riemann. Por muchos años los matemáticos se dieron a la tarea de encontrar las funciones que resultaban de incremetar el número de rectángulos bajo la curva de diversas funciones. No fue sino hasta que Newton y Leibniz acuñaron el teorema fundamental del cálculo que se obtuvieron las herramientas necesarias para estimar \emph{fácilmente} el ára bajo la curva de cualquier función.

\textbf{De aquí en adelante está en construcción, no usar!}

\hypertarget{el-teorema-funamental-del-cuxe1lculo}{%
\subsection{El teorema funamental del cálculo}\label{el-teorema-funamental-del-cuxe1lculo}}

Este teorema establece, en términos muy generales que

\begin{equation}
\int \frac{d}{dx}f = f(x)
\end{equation}

la integral de una derivada es la función misma con la que se obtuvo la derivada. Así, llegamos a la antiderivada.

\hypertarget{antiderivadas-de-funciones}{%
\subsubsection{Antiderivadas de funciones}\label{antiderivadas-de-funciones}}

Bajo la regla de las potencias, que si \(f(x) = ax^n\), \(\int f(x) = \frac{ax^{n+1}}{n+1}\). Si derivamos \(\int f(x)\), veremos que \(f(x) = ax^n\).

Esta sencillísima técnica es efectiva para funciones sin constantes, puesto que la derivada de una constante es 0. Sin embargo debido a que al integrar una función no tenemos ningún conocimiento sobre la existencia o ausencia de una constante (p.~ej. \(y = a + bx\)), es necesario asumir su existencia, de modo que:

\begin{equation}
\int f(x) = \frac{ax^{n+1}}{n+1} + c
\end{equation}

donde \(c\) es la constante de integración.

\hypertarget{aplicaciones-de-la-integraciuxf3n}{%
\subsection{Aplicaciones de la integración}\label{aplicaciones-de-la-integraciuxf3n}}

\hypertarget{cambio-acumulativo}{%
\subsubsection{Cambio acumulativo}\label{cambio-acumulativo}}

El modelo de crecimiento exponencial, basado en la ecuación:

\begin{equation}
    N(t) = e^{rt} \label{eq:ExponN}
\end{equation}

es la base de un gran número de modelos en ecología. Cuando transformamos este modelo en una ecuación diferencial, toma una apariencia un tanto diferente. Comencemos por recordar que si \(f(x) = e^x\), entonces \(f'(x) = e^x\). Si agregamos un parámetro \(r\) como coeficiente de \(x\), tenemos que \(f(x)= e^{rx}\), resultando exactamente en el modelo \eqref{eq:ExponN}.

Para derivar este modelo necesitamos recordar que:

\[\frac{d}{dx}(f \circ g)(x) = (f \circ g)'(x)g'(x)\]
y por analogía \((f \circ g)(x) = e^{rt}\) y \(g(x) = rt\), por lo tanto:

\begin{equation}
\frac{dN}{dt} = r e^{rt}
\end{equation}

donde podemos sustituir \(e^{rt}\) por \(N\), para obtener:

\begin{equation}
\frac{dN}{dt} = r N
\end{equation}

que el cambio inmediato en el número de individuos de la población \(N\), es proporcional a \(rN\). La integral \(\int_{0}^{t} \frac{dN}{dt}\) es entonces equivalente al número de individuos en ese momento \(t\), mientras que \(\int_{0}^{t} N(t)\) corresponde al número total de individuos que hay y ha habido en esa población!

\hypertarget{integraciuxf3n-analuxedtica-vs-numuxe9rica}{%
\subsubsection{Integración analítica vs numérica}\label{integraciuxf3n-analuxedtica-vs-numuxe9rica}}

En matemáticas, el término \emph{analítico} se utiliza para describir el proceso de encontrar con herramientas formales de las matemáticas una solución exacta a un problema. Por ejemplo la antiderivación, aplicando a la inversa las relgas de la derivación es la integración analítica. Hay modelos sin embargo que no se pueden resolver con herramientas analíticas, típicamente, aquellos en que no podemos separar las variables dependientes e independientes. En esos casos hacemos integración \emph{numérica}, que consiste en aproximaciones al valor de la función pero nunca representan las soluciones exactas que se obtienen con integración analítica.

Veamos la ecuación diferencial:

\begin{equation}
\frac{dN}{dt} = rN \label{eq:dif-1}
\end{equation}

corresponde a una función que conocemos como \(N = e^{rt}\), sin embargo podríamos aproximarnos recursivamente a los valores que prediciría \(e^{rt}\) utilizando la definición de la derivada. Sabemos que la derivada es la pendiente de una recta, por lo tanto podríamos comenzar por dibujar una recta con pendiente \(rN\) partiendo desde \(t = 0\), hasta un punto \(0+dx\). El cambio absoluto de \(N\) (\(dN\)) sería entonces \(rN \times dt\), cuyo valorr utilizaremos como punto de partida. Mencioné anteriormente que este método es una aproximación, y esto es porque resulta en un error que se va acumulando entre la función verdadera \(N = e^{rt}\) y la que aproximamos revirtiendo secuencialmente la derivada (figura ).

\begin{figure}

{\centering \includegraphics{Mod-mat-ECO_files/figure-latex/integ-numer-1} 

}

\caption{Diferencias entre integración analítica vs integración numérica, para r = 0.1, dt = 1, y población inicial de N = 1.}\label{fig:integ-numer}
\end{figure}

\hypertarget{unidad-iv-modelos-deterministas-basados-en-ecuaciones-diferenciales-ordinarias}{%
\chapter{Unidad IV: Modelos deterministas basados en ecuaciones diferenciales ordinarias}\label{unidad-iv-modelos-deterministas-basados-en-ecuaciones-diferenciales-ordinarias}}

\hypertarget{modelos-lineales-crecimiento-exponencial-malthusiano}{%
\section{Modelos lineales: Crecimiento exponencial (Malthusiano)}\label{modelos-lineales-crecimiento-exponencial-malthusiano}}

El modelo de crecimiento exponencial también recibe el nombre de ``Malthusiano'', en honor al matemático inglés Thomas Malthus, quien acuñó la teoría demográfica subyacente en la publicación anónima de 1798 \emph{Ensayo sobre el principio de la población}. Este libro trata, no de ecología, sino de las raíces de la pobreza en las sociedades capitalistas, según Malthus, ocasionada por la escasez de recursos. Muchas de las ideas de Malthus, son típicas del pensamiento colonial inglés del período en que vivió, sin embargo, son ecológicamente significativas.

El modelo malthusiano sugiere que, en ausencia de limitaciones, las poblaciones de humanos tenderían a crecer exponencialmente, de ahí que al modelo:

\begin{equation}
 \frac{dN}{dt} = rN \label{eq:Malthus}
\end{equation}

se le llame modelo de crecimiento malthusiano. La intepretación de \(r\) como la tasa de crecimiento sin limitación de alimento, espacio y en ausencia de competencia entre individuos o especies, es la principal aportación de Malthus, y ha servido de base para la formulación de teorías como la evolución.

\hypertarget{interpretaciuxf3n-de-r}{%
\subsection{\texorpdfstring{Interpretación de \(r\)}{Interpretación de r}}\label{interpretaciuxf3n-de-r}}

Más que tasa de crecimiento, \(r\) representa la tasa de cambio del proceso que estamos observando. Estrictamente hablando una población cambia en razón a dos procesos, nacimientos y muertes. De manera relativamente obvia, los nacimientos representan el crecimiento de la población y las muertes del decrecimiento. Por lo tanto, es frecuente considerar que \(r = n - m\), con \(n=\) nacimientos y \(m=\) muertes. Si sustiuimos esta re-interpretación de \(r\), en la ecuación \eqref{eq:Malthus} tenemos:

\begin{equation}
    \frac{dN}{dt} = (n - m)N = nN - mN
\end{equation}

esta re-parametrización del modelo es buena para ejemplificar que si \(n > m \rightarrow r>0\), si \(n < m \rightarrow r < 0\) y que si \(n = m \rightarrow r = 0\). El balance entre nacimientos y muertes entonces es la clave para que una población permanezca constante, crezca o disminuya.

Este balance entre nacimientos y muertes puede también ser concebido como flujos de entrada y salida de un compartimento (figura \ref{fig:nacim-muert}).

\begin{figure}

{\centering \includegraphics[width=12.42in]{Unidad-IV/Malthus-nacim-muert} 

}

\caption{Esquema representativo del modelo de crecimiento malthusiano como flujo de nacimientos y muertes.}\label{fig:nacim-muert}
\end{figure}

\hypertarget{anuxe1lisis-del-modelo}{%
\subsection{Análisis del modelo}\label{anuxe1lisis-del-modelo}}

Hasta el momento hemos visto ya muchas de las propiedades el modelo exponencial:

\begin{enumerate}
\def\labelenumi{\arabic{enumi}.}
\tightlist
\item
  Su derivada es equivalente al valor de la función misma
\item
  Su antiderivada es por lo tanto, igualmente equivalente al valor de la función
\item
  Puede ser representado, como ecuación diferencial, de manera implícita, sin que la variable independiente \(t\) aparezca en el lado derecho de la ecuación \eqref{eq:Malthus}
\item
  Al ser representado de manera implícita, el modelo exponencial es lineal con respecto de la variable dependiente del tiempo
\end{enumerate}

Estas propiedades hacen que analizar el modelo de crecimiento malthusiano sea relativamente sencillo.

\hypertarget{anuxe1lisis-de-estabilidad}{%
\subsection{Análisis de estabilidad}\label{anuxe1lisis-de-estabilidad}}

Este análisis surge de la pregunta ¿bajo qué condiciones una población \(N\) con tasa de crecimiento \(r\) permanece constante? Dada la ecuación \eqref{eq:Malthus} sabemos que si \(\frac{dN}{dt} > 0\), \(N \rightarrow \infty\), por el contrario si \(\frac{dN}{dt} < 0\), \(N \rightarrow 0\). Entonces, para que no haya cambios en \(N\), necesitamos que \(\frac{dN}{dt} = 0\). De esta equivalencia tenemos que si:

\[rN = 0\]
Es debido a que, \(r=0\) si \(N \neq 0\), ó \(N=0\) si \(r \neq 0\). Es decir que el número de individuos \(N\) en una población permanecerá constante sólo bajo dos condiciones, si su tasa de crecimiento es cero, ó no tiene individuos.

\hypertarget{peruxedodos-de-duplicaciuxf3n}{%
\subsection{Períodos de duplicación}\label{peruxedodos-de-duplicaciuxf3n}}

Otra pregunta frecuente en el análisis de poblaciones es ¿cuánto tiempo tardará en duplicarse el tamaño de una población? Para responder esta pregunta, necesitamos la forma cerrada de
la ecuación \eqref{eq:Malthus}:

\begin{equation}
N(t) = N_0 e^{rt} \label{eq:Malthus-c}
\end{equation}

En el análisis de estabilidad encontramos los casos en que \(\frac{dN}{dt} = 0\), por el contrario aquí nos interesa encontrar, con base en la ecuación \eqref{eq:Malthus-c} cuánto tiempo transcurrirá para que \(N(t) = 2N_0\), lo cual podemos conocer resolviendo:

\[2N_0 = N_0 e^{rt}\]
para \(t\):

\begin{align}
e^{rt} &= 2\\ \notag
rt &= \ln 2\\ \notag
t &= \frac{\ln 2}{r}
\end{align}

Es así que podemos representar el período de duplicación de una población cualquiera como una función de su tasa de cambio (figura \ref{fig:periodo-dupl}).

\begin{figure}

{\centering \includegraphics{Mod-mat-ECO_files/figure-latex/periodo-dupl-1} 

}

\caption{Período de duplicación en función de la tasa de crecimiento de una población.}\label{fig:periodo-dupl}
\end{figure}

\hypertarget{estimaciuxf3n-de-la-tasa-de-cambio-r}{%
\subsection{\texorpdfstring{Estimación de la tasa de cambio \(r\)}{Estimación de la tasa de cambio r}}\label{estimaciuxf3n-de-la-tasa-de-cambio-r}}

Para un ecólogo es de suma utilidad conocer la variación de \(r\) para poder medir los efectos del medio ambiente u otras especies sobre los cambios demográficos de una o varias poblaciones. En principio es posible obtener los posibles valores de \(r\) a partir de la ecuación \eqref{eq:Malthus-c}, resolviéndola para \(r\):

\begin{align}
\frac{N}{N_0} &= e^{rt}\\
\ln \left(\frac{N}{N_0} \right) &= rt\\
r &= \frac{\ln N - \ln N_0}{t}
\end{align}

En caso de que tengamos observaciones puntuales de \(N\) y \(N_0\) a lo largo del tiempo.

\hypertarget{implicaciones-del-crecimiento-exponencial}{%
\subsection{Implicaciones del crecimiento exponencial}\label{implicaciones-del-crecimiento-exponencial}}

Una manera de visualizar las repercusiones el crecimiento exponencial y los períodos de duplicación queda de manifiesto con el acertijo:

\emph{Tenemos una población de bacterias que se duplica cada minuto, y tenemos cuatro frascos para albergarla. El primer frasco tardó una hora en llenarse, ¿cuánto tiempo pasará para que se llenen los tres frascos restantes?}

\hypertarget{modelos-no-lineales-crecimiento-loguxedstico-de-una-poblaciuxf3n}{%
\section{Modelos no lineales: Crecimiento logístico de una población}\label{modelos-no-lineales-crecimiento-loguxedstico-de-una-poblaciuxf3n}}

Gracias a Thomas Malthus comprendimos qué sucede con una población cuando llega a estar libre de limitantes externas, sin embargo ¿cómo podemos representar lo que sucede con una población cuando existen dichas limitaciones? La respuesta es el modelo logístico. Este modelo es una variación del modelo malthusiano, pero para regular la tasa neta de crecimiento incorpora una segunda función de \(N\) que al multiplicar a \(r\) se vuelve cero cuando \(N\) alcanza un nivel definido como capacidad de carga (\(K\)).

El crecimiento logístico ocurre en el mundo real más frecuentemente que el malthusiano, pues todos los organismos vivientes crecen a expensas de otros, lo cual también se conoce como denso-dependencia, es decir, que el crecimiento neto de una población disminuye en relación a la cantidad de individuos.

La ecuación del modelo logístico es:

\begin{equation}
    \frac{dN}{dt} = rN \left(1 - \frac{N}{K} \right) \label{eq:logis}
\end{equation}

donde \(K\) es la capacidad de carga del área donde la población con \(N\) individuos vive. Como es evidente la expresión \((1-N/K) \rightarrow 0\) cuando \(N \rightarrow K\). De modo que una población con \(N < K\) tenderá a a crecer. Si expandimos \(rN(1 - N/K) = rN-rN^2/K\) podemos ver que se trata de una parábola. A diferencia del modelo de crecimiento exponencial, donde la tasa de cambio \(r\) es constante, en el modelo logístico, la tasa neta cambia con el tamaño de la población.

\begin{figure}

{\centering \includegraphics{Mod-mat-ECO_files/figure-latex/logis-1} 

}

\caption{Tasa de cambio de N en relación a N en el modelo logístico, con K = 50 y r = 1}\label{fig:logis}
\end{figure}

Algunas propiedades importantes del modelo logístico es que tiene una serie de puntos críticos. Por ejemplo, el tamaño de la población en que ingresa el número máximo de individuos nuevos, el cual corresponde a:

\begin{equation}
\frac{d\left( \frac{dN}{dt} \right)}{dN} = r - 2rN/K = 0
\end{equation}

con lo que tenemos que:

\[ N= \frac{K}{2}\]

El crecimiento máximo se da cuando la población se encuentra al 50\% de la capacidad de carga.

\hypertarget{anuxe1lisis-de-estabilidad-1}{%
\subsubsection{Análisis de estabilidad}\label{anuxe1lisis-de-estabilidad-1}}

Al igual que con el crecimiento malthusiano, nos interesa encontrar en qué situaciones la población será constante, por lo que tenemos que resolver:

\[\frac{dN}{dt} = 0\]
para encontrar que hay dos puntos de equilibrio, \(N = 0\) y \(N = K\). Estos puntos tienen propiedades diferentes. Por ejemplo, ¿qué sucede con el \(N\) si a \(N=0\) le sumamos una cantidad muy pequeña \(\varepsilon\)? ¿\(N\) tenderá a regresar a \(0\)? Un ejemplo de este tipo de dinámicas en el mundo real se presenta con un péndulo, el cual tiene dos puntos de equilibrio, uno estable (donde el péndulo descansa en ausencia de una fuerza lateral), y otro inestable

\begin{figure}

{\centering \includegraphics[width=10.22in]{Unidad-IV/Puntos-equilibrio} 

}

\caption{Puntos de equilibrio de un péndulo.}\label{fig:equilibrio}
\end{figure}

Las mismas preguntas aplican para \(N=K\). Para responderlas tenemos que hacer un paréntesis y echar un vistazo a la integración numérica.

\hypertarget{integraciuxf3n-numuxe9rica}{%
\paragraph{Integración numérica}\label{integraciuxf3n-numuxe9rica}}

Esta técnica de integración está basada en la definición de la derivada como límite:

\begin{equation}
f'(x) = \lim \limits_{h \rightarrow 0} \frac{f(x + h) - f(x)}{h} \label{eq:deriv-Euler}
\end{equation}

donde se asume que \(h = dx\), con \(dx \approx 0\), es decir, que \(dx\) es una cantidad arbitrariamente pequeña. De modo que sustituimos \(dx\) por \(h\) en \eqref{eq:deriv-Euler} y resolvemos para \(f(x + dx)\) para obtener:

\begin{equation}
f(x + dx) = f(x) + f'(x) \cdot dx \label{eq:Euler}
\end{equation}

Este método es el más sencillo de los métodos numéricos de integración, y ha servido de punto de partida para el desarrollo de métodos más complejos y precisos.

Para integrar numéricamente el modelo logístico, utilizando la técnica de Euler \eqref{eq:Euler}, para lo que necesitamos conocer:

\begin{enumerate}
\def\labelenumi{\arabic{enumi}.}
\tightlist
\item
  \(r = 0.5\): la tasa neta de crecimiento poblacional
\item
  \(K = 50\): la capacidad de carga poblacional
\item
  \(N_0 = 1\): El tamaño inicial de la población
\item
  \(dt = 1\): El tamaño de los incrementos de la variable independiente \(t\)
\item
  \(t_{max}\): El tiempo máximo de simulación
\end{enumerate}

Estos valores lo utilizaremos para conocer \(dN/dt\):

\[\frac{dN_0}{dt} = 0.5 \times N_0 \times \left (1 - \frac{N_0}{50} \right)=0.49\]

Por lo tanto \(N_1\) es:

\[N(1) = N_0 + \frac{dN_0}{dt} \times dt = 1 + 0.49 \times 1 = 1.49\]
Este proceso se repite hasta que \(t = t_{max}\), y obtengamos \(N(t_{max})\). El código para integrar esta ecuación en \textbf{R} es:

\begin{Shaded}
\begin{Highlighting}[]
\CommentTok{\#Parámetros}
\NormalTok{dt }\OtherTok{\textless{}{-}} \DecValTok{1}
\NormalTok{r }\OtherTok{\textless{}{-}} \FloatTok{0.5}
\NormalTok{K }\OtherTok{\textless{}{-}} \DecValTok{50}

\CommentTok{\# Variables}
\NormalTok{N }\OtherTok{\textless{}{-}} \FunctionTok{numeric}\NormalTok{(}\DecValTok{25}\NormalTok{)}
\NormalTok{dN }\OtherTok{\textless{}{-}} \FunctionTok{numeric}\NormalTok{(}\DecValTok{25}\NormalTok{)}

\CommentTok{\# Valores iniciales}
\NormalTok{N[}\DecValTok{1}\NormalTok{] }\OtherTok{\textless{}{-}} \DecValTok{1}
\NormalTok{dN[}\DecValTok{1}\NormalTok{] }\OtherTok{\textless{}{-}}\NormalTok{ r}\SpecialCharTok{*}\NormalTok{N[}\DecValTok{1}\NormalTok{] }\SpecialCharTok{*}\NormalTok{ (}\DecValTok{1} \SpecialCharTok{{-}}\NormalTok{ N[}\DecValTok{1}\NormalTok{]}\SpecialCharTok{/}\NormalTok{K)}

\CommentTok{\# Integración}
\ControlFlowTok{for}\NormalTok{(t }\ControlFlowTok{in} \DecValTok{2}\SpecialCharTok{:}\DecValTok{25}\NormalTok{)\{}
\NormalTok{    N[t] }\OtherTok{\textless{}{-}}\NormalTok{ N[t}\DecValTok{{-}1}\NormalTok{] }\SpecialCharTok{+}\NormalTok{ dN[t}\DecValTok{{-}1}\NormalTok{] }\SpecialCharTok{*}\NormalTok{ dt}
\NormalTok{    dN[t] }\OtherTok{\textless{}{-}}\NormalTok{ r}\SpecialCharTok{*}\NormalTok{N[t] }\SpecialCharTok{*}\NormalTok{ (}\DecValTok{1} \SpecialCharTok{{-}}\NormalTok{ N[t]}\SpecialCharTok{/}\NormalTok{K)}
\NormalTok{\}}
\end{Highlighting}
\end{Shaded}

\begin{figure}

{\centering \includegraphics{Mod-mat-ECO_files/figure-latex/logis-Euler-1} 

}

\caption{Simulación del modelo logístico con el método de Euler, con el código descrito aquí.}\label{fig:logis-Euler}
\end{figure}

\hypertarget{anuxe1lisis-de-estabilidad-2}{%
\paragraph{Análisis de estabilidad}\label{anuxe1lisis-de-estabilidad-2}}

Aquí veremos si los puntos de equilibrio \(N=0\) y \(N=K\) son estables o inestables. Para ello simularemos la ecuación con ambas condiciones iniciales, y les sumaremos una cantidad pequeña \(\varepsilon = 0.01\) para ver qué sucede con \(N\) en el resto de la simulación. En analogía con el péndulo, esperaríamos que \(N \rightarrow N^*\), el punto de equilibrio estable (\(0\) ó \(K\)) al aplicar la pertubación \(\varepsilon\) con ambas condiciones de inicio.

\begin{Shaded}
\begin{Highlighting}[]
\NormalTok{dt }\OtherTok{\textless{}{-}} \DecValTok{1}
\NormalTok{r }\OtherTok{\textless{}{-}} \FloatTok{0.5}
\NormalTok{K }\OtherTok{\textless{}{-}} \DecValTok{50}

\CommentTok{\# Variables}
\NormalTok{N }\OtherTok{\textless{}{-}} \FunctionTok{numeric}\NormalTok{(}\DecValTok{25}\NormalTok{)}
\NormalTok{dN }\OtherTok{\textless{}{-}} \FunctionTok{numeric}\NormalTok{(}\DecValTok{25}\NormalTok{)}

\NormalTok{E }\OtherTok{\textless{}{-}} \FloatTok{0.01}

\CommentTok{\# Valores iniciales}
\NormalTok{N[}\DecValTok{1}\NormalTok{] }\OtherTok{\textless{}{-}} \DecValTok{0} \SpecialCharTok{+} \FloatTok{0.01}
\NormalTok{dN[}\DecValTok{1}\NormalTok{] }\OtherTok{\textless{}{-}}\NormalTok{ r}\SpecialCharTok{*}\NormalTok{N[}\DecValTok{1}\NormalTok{] }\SpecialCharTok{*}\NormalTok{ (}\DecValTok{1} \SpecialCharTok{{-}}\NormalTok{ N[}\DecValTok{1}\NormalTok{]}\SpecialCharTok{/}\NormalTok{K)}

\CommentTok{\# Integración}
\ControlFlowTok{for}\NormalTok{(t }\ControlFlowTok{in} \DecValTok{2}\SpecialCharTok{:}\DecValTok{25}\NormalTok{)\{}
\NormalTok{    N[t] }\OtherTok{\textless{}{-}}\NormalTok{ N[t}\DecValTok{{-}1}\NormalTok{] }\SpecialCharTok{+}\NormalTok{ dN[t}\DecValTok{{-}1}\NormalTok{] }\SpecialCharTok{*}\NormalTok{ dt}
\NormalTok{    dN[t] }\OtherTok{\textless{}{-}}\NormalTok{ r}\SpecialCharTok{*}\NormalTok{N[t] }\SpecialCharTok{*}\NormalTok{ (}\DecValTok{1} \SpecialCharTok{{-}}\NormalTok{ N[t]}\SpecialCharTok{/}\NormalTok{K)}
\NormalTok{\}}

\CommentTok{\#Almacenamos el resultado con la condición inicial N=0+E en otro objeto N0}
\NormalTok{N0 }\OtherTok{\textless{}{-}}\NormalTok{ N}

\CommentTok{\#Repetimos para N=K+E}
\NormalTok{N }\OtherTok{\textless{}{-}} \FunctionTok{numeric}\NormalTok{(}\DecValTok{25}\NormalTok{)}
\NormalTok{dN }\OtherTok{\textless{}{-}} \FunctionTok{numeric}\NormalTok{(}\DecValTok{25}\NormalTok{)}

\NormalTok{N[}\DecValTok{1}\NormalTok{] }\OtherTok{\textless{}{-}}\NormalTok{ K }\SpecialCharTok{+} \FloatTok{0.01}
\NormalTok{dN[}\DecValTok{1}\NormalTok{] }\OtherTok{\textless{}{-}}\NormalTok{ r}\SpecialCharTok{*}\NormalTok{N[}\DecValTok{1}\NormalTok{] }\SpecialCharTok{*}\NormalTok{ (}\DecValTok{1} \SpecialCharTok{{-}}\NormalTok{ N[}\DecValTok{1}\NormalTok{]}\SpecialCharTok{/}\NormalTok{K)}

\CommentTok{\# Integración}
\ControlFlowTok{for}\NormalTok{(t }\ControlFlowTok{in} \DecValTok{2}\SpecialCharTok{:}\DecValTok{25}\NormalTok{)\{}
\NormalTok{    N[t] }\OtherTok{\textless{}{-}}\NormalTok{ N[t}\DecValTok{{-}1}\NormalTok{] }\SpecialCharTok{+}\NormalTok{ dN[t}\DecValTok{{-}1}\NormalTok{] }\SpecialCharTok{*}\NormalTok{ dt}
\NormalTok{    dN[t] }\OtherTok{\textless{}{-}}\NormalTok{ r}\SpecialCharTok{*}\NormalTok{N[t] }\SpecialCharTok{*}\NormalTok{ (}\DecValTok{1} \SpecialCharTok{{-}}\NormalTok{ N[t]}\SpecialCharTok{/}\NormalTok{K)}
\NormalTok{\}}

\NormalTok{NK }\OtherTok{\textless{}{-}}\NormalTok{ N}
\end{Highlighting}
\end{Shaded}

\begin{figure}

{\centering \includegraphics{Mod-mat-ECO_files/figure-latex/Estabilidad-1} 

}

\caption{Simulaciones para verificar tipo de puntos de equilibrio}\label{fig:Estabilidad}
\end{figure}

\hypertarget{modelos-en-tiempo-contuxednuo-vs-discreto}{%
\section{Modelos en tiempo contínuo vs discreto}\label{modelos-en-tiempo-contuxednuo-vs-discreto}}

Todos los modelos que hemos visto en la presente unidad son dinámicos, representan procesos que cambian con el tiempo y su principal insumo son la condiciones presentes del sistema, las cuales determinan el estado futuro. Aprendimos a integrar ecuaciones diferenciales utilizando la definición de la derivada, y revirtiéndola asumiendo que \(h \approx dt\). Este proceso, es en esencia una discretización del tiempo, es decir, en lugar de evaluar \(f(t)\) para todos los puntos \(t\) posibles, lo aproximamos para intervalos separados por \(dt\). Por lo tanto las ecuaciones diferenciales integradas con métodos numéricos tienen algunas similitudes con los modelos en tiempo discreto:

Los modelos en tiempo discreto se conocen como ecuaciones en diferencias finitas, en comparación con las ecuaciones diferenciales que representan el tiempo de manera contínua.

La difefencia entre estos modelos consiste en los tiempos en que ocurren los eventos que representan. En los modelos de crecimiento poblacional los eventos representan la incorporación de invidiuos nuevos a la población. Por lo tanto, en un modelo de crecimiento contínuo los individuos nuevos entrar a la población ``todo el tiempo''.

\hypertarget{crecimiento-exponencial-en-tiempo-discreto}{%
\subsection{Crecimiento exponencial en tiempo discreto}\label{crecimiento-exponencial-en-tiempo-discreto}}

Las ecuaciones de los modelos contínuo y discreto corresponden a:

\begin{align}
\frac{dN}{dt} & = rN \label{eq:exp-cont} \\
N_{t+1} & = \lambda N \label{eq:exp-disc} 
\end{align}

donde las tasas de cambio son \(r\) y \(\lambda\). En el caso del modelo en tiempo discreto \eqref{eq:exp-disc}, se utiliza el parámetro \(\lambda\), puesto que \(\lambda = e^r\)

El crecimiento discreto tiene sentido biológico cuando los nacimientos de las poblaciones de algún organismo se sincronizan en un período de tiempo muy corto. El crecimiento exponencial contínuo, en contraste, representa las dinámicas de organismos que pueden reproducirse en cualquier momento. En un modelo de crecimiento discreto, los individuos nuevos entran a la población solamente en un momento preciso. La cantidad de individuos que ingresa a la pobalación al final de cada intervalo es la misma que en un modelo en tiempo contínuo (Figura \ref{fig:disc-cont}). Por lo tanto, la selección de modelo tiene que estar informada por la biología del problema.

\begin{figure}

{\centering \includegraphics{Mod-mat-ECO_files/figure-latex/disc-cont-1} 

}

\caption{Diferencia entre modelos en tiempo contínuo (línea discontínua roja) y discreto (escalones en azul).}\label{fig:disc-cont}
\end{figure}

La escala exponencial del parámetro \(\lambda\) hace que la solución sea exacta, en comparación con la aproximación numérica con la ecuación de Euler del modelo en tiempo contínuo.

\hypertarget{crecimiento-loguxedstico-en-tiempo-discreto}{%
\subsection{Crecimiento logístico en tiempo discreto}\label{crecimiento-loguxedstico-en-tiempo-discreto}}

El modelo discreto es:

\begin{equation}
N_{t+1} = \lambda N \left (1 - \frac{N}{K} \right)
\end{equation}

en este caso \(\lamda\) está en la misma escala que en la ecuación \eqref{eq:exp-disc}. Como ya vimos anteriormente, esta versión del modelo de crecimiento logístico, tiene una serie de propiedades distintas al modelo contínuo, por lo que nos enfocaremos en este aspecto. Para ello, transformaremos sustuiremos \(N\) por \(x_t = N_t/K\), para analizar el tamaño de \(N\) en relación a la capacidad de carga:

\begin{equation}
x_{t+1} = b(1-x_t)x_t \label{eq:logis-disc}
\end{equation}

El análisis del comportamiento del modelo logístico es más fácil si se hace gráficamente, en especial porque puede ser complejo, porque puede ser caótico (Figura \ref{fig:caos-logis})

\begin{figure}

{\centering \includegraphics{Mod-mat-ECO_files/figure-latex/caos-logis-1} 

}

\caption{Ejemplos del comportamiento dinámico del modelo de crecimiento logístico, para tres valores de la tasa de crecimiento. En las figuras de la izquiera se muestra la correspondencia entre el tamaño de la población en t y el tamaño en t+1.}\label{fig:caos-logis}
\end{figure}

\hypertarget{implicaciones-para-sistemas-ecoluxf3gicos}{%
\subsubsection{Implicaciones para sistemas ecológicos}\label{implicaciones-para-sistemas-ecoluxf3gicos}}

Las dinámicas producidas por el sencillísimo modelo logístico en tiempo discreto sugieren que un sistema ecológico con esas características bien podría tener un comportamiento caótico. Como ecólogos es preciso preguntarnos si sólo se trata de artefactos matemáticos producto de la simplificación, o si en posible su ocurrencia. Modelos más complejos que toman en cuenta la existencia de diferentes sexos, por ejemplo, tienen un comportamiento más estable que los modelos sin esas estructuras \citep[\citet{Scheuring2002IsCD}]{berryman1989ecological}.

La diferencia esencial entre el caos producido por el modelo logístico y el comportamiento errático de muchos sistemas ecológicos está en que el caos forma parte del funcionamiento del sistema, es determinístico, mientras que las variaciones de un sistema ecológico generalmente se consideran como fuerzas externas que cambian con el tiempo y que afectan al sistema.

\hypertarget{unidad-v-introducciuxf3n-al-uxe1lgebra-de-matrices}{%
\chapter{Unidad V: Introducción al álgebra de matrices}\label{unidad-v-introducciuxf3n-al-uxe1lgebra-de-matrices}}

\hypertarget{sistemas-lineales}{%
\section{Sistemas lineales}\label{sistemas-lineales}}

Un sistema lineal es un conjunto de ecuaciones lineales donde hay, idealmente, tantas variables como ecuaciones y viceversa. Cuando esta condición se cumple, se dice que el sistema está \textbf{completo}. La contraparte de estos sistemas son \textbf{incompletos}, y pueden tener más variables que ecuaciones ó más ecuaciones que variables.

Un sistema completo de ecuaciones lineales con dos variables sería:

\[2x + 3y = 4\]

\[4x - 2y =10\]
Generalizando tenemos que un cualquier sistema lineal puede ser:

\begin{align}
a_{1,1}x_1+a_{2,1}x_2 + \dots + a_{n,1}x_n & = b_1\\
a_{1,2}x_1+a_{2,2}x_2 + \dots + a_{n,2}x_n & = b_2\\
\vdots & \\
a_{1, n}x_1+a_{2,n}x_2 + \dots + a_{n,n}x_n & = b_n
\end{align}

Cuando nos enfrentamos a un sistema como estos, la meta es encontrar los valores de \(x_1, x_2, \dots , x_3\) que satisfacen las igualdades. Para sistemas de dos ecuaciones, es fácil econtrar la solución utilizando una variedad de métodos, el primero de ellos, por graficación.

\hypertarget{soluciuxf3n-gruxe1fica}{%
\subsection{Solución gráfica}\label{soluciuxf3n-gruxe1fica}}

La solución para el sistema de ecuaciones del primer ejemplo:

\begin{align}
2x + 3y & = 4\\
4x - 2y & =10
\end{align}

sería:

\begin{figure}

{\centering \includegraphics[width=41.67in]{Unidad-V/Soln-grafica-1} 

}

\caption{Solución gráfica del sistema de ecuaciones para dos variables}\label{fig:unnamed-chunk-41}
\end{figure}

\hypertarget{soluciuxf3n-algebruxe1ica}{%
\subsection{Solución algebráica}\label{soluciuxf3n-algebruxe1ica}}

El método más sencillo en esta ocasión sería resolver la ecuación 2 para \(y\), pues la podemos simplificar dividiendo ambos lados entre 2:

\[ (4x - 2y)/2 = (10)/2 \rightarrow 2x - y = 5\]
de donde obtenemos que:

\[y = 2x-5\]

Este valor de \(y\) en térmonos de \(x\), lo podemos sustituir en, la ecuación 1:

\[2x+3(2x-5) = 4\]

y resolvemos para \(x\) para obtener:

\[8x = 19 \rightarrow x = 19/8=2.375\]

Y entonces sustituimos a \(x\) en:

\(y = 2(2.375) - 5 = -0.25\)

Como se puede apreciar, resolver simples ecuaciones lineales puede ser laborioso. Es por ello que se desarrolló el álgebra de matrices, para resolver de eficientemente, sistemas de ecuaciones lineales con muchas variables. En esta unidad revisaremos brevemente los conceptos y técnicas más importantes del álgebra de matrices.

\hypertarget{sistemas-lineales-sin-soluciuxf3n}{%
\subsection{Sistemas lineales sin solución}\label{sistemas-lineales-sin-soluciuxf3n}}

En ocasiones, aunque un sistema de ecuaciones sea aparentemente completo, puede no tener solución. Por ejemplo el sistema:

\begin{align}
x + y & = 1 \\
2x + 2y & = 4
\end{align}

no tiene solución, lo cual se manifiesta gráficamente como dos líneas paralelas. Esto quiere decir que no existe un par de valores \(x\) y \(y\) que satisfagan las ecuaciones del sistema.

\begin{figure}

{\centering \includegraphics[width=41.67in]{Unidad-V/Sistema-no-solucion} 

}

\caption{Representación gráfica de un sistema linea sin solución.}\label{fig:sin-soln}
\end{figure}

Si intentamos sumar o restar ambas ecuaciones terminamos con una contradicción. Tomemos por ejemplo la primera ecuación, multiplicada por \(-2\) para obtener:

\[(x + y = 1) \times -2 \rightarrow -2x-2y = -2\]

Y al sumarla a la segunda ecuación:

\[2x - 2x + 2y - 2y = 4-2 \rightarrow 0 = 2\]

\hypertarget{resolviendo-sistemas-lineales-muxe1s-complejos}{%
\subsection{Resolviendo sistemas lineales más complejos}\label{resolviendo-sistemas-lineales-muxe1s-complejos}}

Cualquier sistema de ecuaciones lineales son solución se puede revolver con los métodos que hemos aprendido hasta el momento. Como hemos visto, en estos métodos es fácil cometer errores, máxime con sistemas más complejos con muchas variables y ecuaciones. El método ideal sería aquel que:

\begin{enumerate}
\def\labelenumi{\arabic{enumi}.}
\tightlist
\item
  Evita manipulaciones riesgosas de las ecuaciones
\item
  Permite encontrar todas las soluciones posibles
\item
  Permite identificar sistemas sin solución
\end{enumerate}

En la tarea de aprender a utilizar dichos métodos cabe resaltar que las incógnitas (\(x, y, z, \dots\)) juegan un papel secundario, lo más importante son los coeficientes de cada una en las distintas ecuaciones.

Para comenzar a visualizar un sistema lineal como matriz, echemos un vistazo al siguiente sistema:

\begin{align}
x+y+z+t & = 1 \\
x-y-z+t & = 3 \\
-x-y+z-t & = 1\\
-3x+y-3z-3t &= 4
\end{align}

cuyos coeficientes podemos representar en una matriz como:

\[A = \left[
\begin{array}{rrrr}
1 & 1 & 1 & 1 \\
1 &-1 &-1 & 1 \\
-1&-1 & 1 &-1 \\
-3& 1 &-3 &-3
\end{array}
\right]\]

Nota que cada columna corresponde con cada una de las variables \(x, y, z, t\), y los números que contiene la matriz \(A\) son sus coeficientes. Usando este mismo arreglo, podemos representar todo el sistema de ecuaciones como el producto de un vector \(x, y, z, t\) y la matriz \(A\):

\begin{equation}
\left[
\begin{array} {r}
x \\ y \\ z \\ t
\end{array}
\right]
\left[
\begin{array}{rrrr}
1 & 1 & 1 & 1 \\
1 &-1 &-1 & 1 \\
-1&-1 & 1 &-1 \\
-3& 1 &-3 &-3
\end{array} 
\right] = 
\left[\begin{array}{r}
1 \\ 3 \\ 1 \\ 4
\end{array}\right] \label{eq:sist-matriz}
\end{equation}

Dado que para las operaciones necesarias para resolverlo, sólo necesitamos la matriz de coeficientes \(A\) y el vector de equivalencias podemos representar el sistema con su matriz aumentada:

\begin{equation}
\left[
\begin{array}{rrrr|l}
1 & 1 & 1 & 1 & 1\\
1 &-1 &-1 & 1 & 3\\
-1&-1 & 1 &-1 & 1\\
-3& 1 &-3 &-3 & 4
\end{array} \right] \label{eq:Aumentada}
\end{equation}

Para resolver sistemas como este necesitamos aprender a hacer algunas operaciones con filas que no alteran el estado del sistema lineal:

\begin{enumerate}
\def\labelenumi{\arabic{enumi}.}
\tightlist
\item
  Intercambiar filas
\item
  Multiplicar una fila por un escalar \(c \neq 0\)
\item
  Sumar dos filas
\end{enumerate}

Ninguna de estas operaciones afectan las soluciones posibles del sistema (si las hay). Utilizando estas operaciones, la meta es transformar la matriz aumentada (\eqref{eq:Aumentada}), en un a matriz que contenga ceros bajo la diagonal (\textbf{Triángulo superior}), y así poder sustituir los valores comenzando con la última fila. El sistema de ecuaciones \eqref{eq:sist-matriz}, no tiene solución, por lo que solamente nos sirve, por el momento para establecer los conceptos. A continuación lo aplicamos con un sistema más sencillo.

\hypertarget{ejemplo}{%
\subsection{Ejemplo}\label{ejemplo}}

Dado el sitema linea:

\begin{align}
2x - 1y + 3z &= 9\\
3x + 5y - z  &= 10\\
4x + 2y - 3z &= -1
\end{align}

la matriz aumentada es:

\begin{equation}
\left[
\begin{array}{rrr|l}
 2 & -1& 3 & 9 \\
 3 & 5 &-1 & 10\\
 4 & 2 &-3 &-1
\end{array}
\right]
\begin{array}
\ (F_1) \\ (F_2) \\ (F_3)
\end{array}
\end{equation}

A continuación se muestra en el lado izquierdo, las operaciones correspondientes a las filas de las matrices:

\begin{equation}
\begin{array}{r}
(F_1) \\ -3(F_1)+2(F_2) \\ 2(F_1)-(F_3)
\end{array}
\left[
\begin{array}{rrr|r}
2 &-1 & 3 & 9 \\
0 &13 &-11&-7 \\
0 &-4 & 9 &19
\end{array}
\right]
\begin{array}{r}
\ (F_4) \\ (F_5) \\ (F_6)
\end{array}
\end{equation}

Ahora la fila 1 se llama \((F_4)\), hicimos \(-3(F_1)+2(F_2)\) para eliminar el coeficiente \(a_{2,1} = 3\), y ahora lo llamamos \((F_5)\). Lo único que debemos hacer ahora, es encontrar cómo eliminar \(a_{2,3}\) en \((F_6)\), para obtener la matriz con ceros bajo la diagonal:

\begin{equation}
\begin{array}{r}
(F_4) \\ (F_5) \\ 4(F_5)+13(F_6)
\end{array}
\left[
\begin{array}{rrr|r}
2 &-1 & 3 & 9 \\
0 &13 &-11&-7 \\
0 & 0 &73 &219
\end{array}
\right]
\begin{array}{r}
(F_7) \\ (F_8) \\ (F_9)
\end{array}
\end{equation}

A partir de aquí, tenemos una ecuación on una sola variable:

\[73z = 219\]
Con la cual podemos sustituir desde \((R_9)\) hacia \((R_7)\) para obtener los valores de \(x\) y \(y\) que satisfacen al sistema lineal.

\hypertarget{sistemas-indeterminados}{%
\subsection{Sistemas indeterminados}\label{sistemas-indeterminados}}

Cuando hay más variables que ecuaciones, no es posible encontrar una solución única al sistema, veamos un ejemplo de estos sistemas:

\begin{align}
2x+2y-z &= 1 \\
2x-y+z &= 2
\end{align}

cuya matriz aumentada es:

\begin{equation}
\left[
\begin{array}{rrr|r}
2 & 2 & -1& 1 \\
2 & -1& 1 & 2
\end{array} \right]
\begin{array}
(F_1) \\ (F_2)
\end{array}
\end{equation}

Lo primero que notamos es que, al haber más variables que ecuaciones, la matriz de coeficientes, tiene más columnas que filas. Entonces, vamos a comenzar por eliminar el coeficiente \(a_{2,1} = 2\), haciendo \((F_1)-F_2\):

\begin{equation}
\begin{array}{r}
(F_1) \\ (F_1) - (F_2)
\end{array}
\left[ 
\begin{array}{rrr|r}
2 & 2 &-1 & 1\\
0 & 3 &-2 &-1
\end{array}
\right]
\begin{array}{r}
(F_3) \\ (F_4)
\end{array}
\end{equation}

A partir de este punto sólo podemos expresar las soluciones para \(x\) y \(y\) en términos de \(z\):

\begin{align}
2x + 2y - z &= 1 \\
3y - 2z & = -1
\end{align}

Por lo tanto:

\[y = - \frac{1}{3} + \frac{2}{3}z\]

y sustituyendo \(y\) en la primera ecuación encontramos que:

\[x = \frac{5}{6} - \frac{1}{6}z\]

Para reportar un conjunto de soluciones para todo el sistema, podemos crear una variable ficticia \(t = z\), con lo que la solución al sistema es:

\[\left\{(x, y, z): x = \frac{5}{6} - \frac{1}{6}z, y = - \frac{1}{3} + \frac{2}{3}z, z = t \right\}\]

\hypertarget{operaciones-elementales-con-matrices}{%
\section{Operaciones elementales con matrices}\label{operaciones-elementales-con-matrices}}

Las operaciones elementales con matrices son similares a las aritméticas tradicionales. Para ejemplificarlas aquí denotamos a una matriz cualquiera con letras mayúsculas.

\hypertarget{suma-de-matrices}{%
\subsection{Suma de matrices}\label{suma-de-matrices}}

Si las matrices \(A\) y \(B\) tienenlas mismas dimensiones \(m \times n\) (\(m\) = número de filas, \(n\) = número de columnas):

\[A + B = C\]
La adición de \(A\) con \(B\) resulta en una tercera matriz \(C\) con dimensiones \(m \times n\), tanto si \(m = n\) como \(m \neq n\), es decir si las matrices \(A\) y \(B\) son cuadradas (\(m = n\)) como rectangulares.

Para esto tenemos que considerar qué sucede con los elementos de \(A\) y \(B\) dentro de \(C\):

\begin{equation}
A = \left[ 
\begin{array}{ccc}
a_{1, 2} & \dots & a_{1, n} \\
\vdots   & \ddots & \vdots \\
a_{n, 1} & \dots & a_{n, n}
\end{array} \right]
\end{equation}

\begin{equation}
B = \left[ 
\begin{array}{ccc}
b_{1, 2} & \dots & b_{1, n} \\
\vdots   & \ddots& \vdots \\
b_{n, 1} & \dots & b_{n, n}
\end{array} \right]
\end{equation}

\begin{equation}
C = \left[ 
\begin{array}{ccc}
a_{1, 2} + b_{1, 2} & \dots & a_{1, n} + b_{1, n} \\
\vdots   & \ddots & \vdots \\
a_{n, 1} + b_{n, 1} & \dots & a_{n, n} + b_{n, n}
\end{array} \right]
\end{equation}

Es claro que también:

\[B+A = C\]

Lo cual resulta importante resaltar pues la conmutatividad no funciona para la multiplicación de matrices. Otras cualidades de la adición de matrices, dadas las condiciones de dimensiones son:

\begin{enumerate}
\def\labelenumi{\arabic{enumi}.}
\tightlist
\item
  \((A+B)+C=A+(B+C)\)
\item
  \(A + 0 = A\)
\end{enumerate}

\hypertarget{multiplicaciuxf3n-por-un-escalar}{%
\subsection{Multiplicación por un escalar}\label{multiplicaciuxf3n-por-un-escalar}}

Es frecuente multiplicar una matriz por una constante \(c\), tal que \(c \times A = cA\):

\begin{equation}
c \times A = \left[ 
\begin{array}{ccc}
c  \cdot  a_{1, 2} & \dots & c \cdot a_{1, n} \\
\vdots   & \ddots & \vdots \\
c \cdot  a_{n, 1} & \dots & c  \cdot  a_{n, n}
\end{array} \right]
\end{equation}

\hypertarget{transposiciuxf3n}{%
\subsection{Transposición}\label{transposiciuxf3n}}

La transposición es el cambio de filas y columnas por columnas y filas, por ejemplo:

\begin{equation}
A = \left[ 
\begin{array}{rrr}
1 & 2 & 3 \\
4 & 5 & 6
\end{array}
\right]
\end{equation}

Con lo que la matriz \(A\) transpuesta \(A^T = A'\) es:

\begin{equation}
A' = \left[
\begin{array}{rr}
1 & 4 \\
2 & 5 \\
3 & 6
\end{array}
\right]
\end{equation}

\hypertarget{multiplicaciuxf3n-de-matrices}{%
\subsection{Multiplicación de matrices}\label{multiplicaciuxf3n-de-matrices}}

Para entender la multiplicación de matrices es necesario definirla, pues es más complicada que cualquiera de las operaciones que hemos visto anteriormente.

Supongamos que \(C\) es una matriz producto de la multiplicación de las matrices \(A\) y \(B\):

\[C = AB\]

Para que esto sea posible es necesario que \(A\) y \(B\) tengan dimensiones específicas pues:

\begin{equation}
C = \left[
\begin{array}{rrr}
 &        & \\
 & c_{ij} & \\
 &        &
\end{array}
\right] =
\left[ 
\begin{array}{ccc}
 & & & \\
a_{i, 1} & \dots & a_{i, l} \\
 & & &
\end{array}
\right] \times
\left[ 
\begin{array}{ccc}
 & b_{1, j} & \\
 & \vdots & \\
 & b_{l, j} & 
\end{array}
\right] \label{eq:mat-mult}
\end{equation}

Como resulta evidente, el número de \textbf{columnas} de la matriz \(A\) debe ser igual al número de \textbf{filas} de la matriz \(B\). De modo que podemos resumir que las dimensiones deben ser:

\begin{itemize}
\tightlist
\item
  \(A \rightarrow m \times l\)
\item
  \(B \rightarrow l \times n\)
\end{itemize}

\hypertarget{ejemplos-3}{%
\subsubsection{Ejemplos}\label{ejemplos-3}}

\begin{equation}
A = \left[ 
\begin{array}
\ 1 & 2 & 3 \\
-1& 0 & 4
\end{array}
\right]
\end{equation}

y

\begin{equation}
B = \left[ 
\begin{array}
\ 1 & 2 & 3 & -3\\
0 &-1 & 4 & 0 \\
-1& 0 &-2 & 1
\end{array}
\right]
\end{equation}

Solución:

\begin{equation}
C = AB = \left[ 
\begin{array}
\ 1 & 2 & 3 \\
-1& 0 & 4
\end{array}
\right]

\left[ 
\begin{array}
\ 1 & 2 & 3 & -3\\
0 &-1 & 4 & 0 \\
-1& 0 &-2 & 1
\end{array}
\right]
\end{equation}

Como vimos en la definición de multiplicación, tenemos que muiltiplicar las filas de \(A\) por las columnas de \(B\), para así obtener los elementos de \(C\). Hagamos entonces las operaciones elemento por elemento:

\begin{equation}
c_{1, 1} = \left[ 
\begin{array}{rrr}
\ 1 & 2 & 3
\end{array}
\right]
\left[\begin{array}
\ 1 \\ 0 \\ -1 \end{array}
\right] = (1)(1) + (2)(0) + (3)(-1)= -2
\end{equation}

\begin{equation}
c_{1, 2} = \left[
\begin{array}
1 & 2 & 3
\end{array}
\right]
\left[ 
\begin{array}
\ 2 \\ -1 \\0
\end{array}
\right] = (1)(2) + (2)(-1) + (3)(0) = 0
\end{equation}

Y así sucesivamente para obtener:

\begin{equation}
C = \left[ 
\begin{array}
\ c_{1, 1} & c_{1, 2} & c_{1, 3} & c_{1, 4} \\
c_{2, 1} & c_{2, 2} & c_{2, 3} & c_{2, 4} 
\end{array}
\right] =
\left[ 
\begin{array}
\ -2 & 0 & 5 & 0 \\
-5 & -2 & -11 & 7
\end{array}
\right]
\end{equation}

El resultado es una matriz con tantas filas como \(A\) y tantas columnas como \(B\). Es decir que si las dimensiones de \(A\) son \(m \times l\), y las de \(B\) son \(l \times n\), \(C\) tendrá dimensiones \(m \times n\).

\hypertarget{ejercicio-2}{%
\paragraph{Ejercicio}\label{ejercicio-2}}

Haz a mano los cálculos de \(c_{i, j}\) que nos faltaron para obtener los valores de \(C\).

\hypertarget{determinantes-y-matrices-inversas}{%
\section{Determinantes y matrices inversas}\label{determinantes-y-matrices-inversas}}

\hypertarget{matriz-identidad}{%
\subsection{Matriz identidad}\label{matriz-identidad}}

Antes de continuar, consideremos la posible existencia de una matriz \(I\) tal que \(A I = A\), es decir, una matriz que multiplicada por otra matriz de como resultado la matriz original. Esta matriz se conoce como la matriz identidad, y consiste de una matriz cuadrada con ceros en todas las posiciones excepto en la diagonal, por ejemplo:

\begin{equation}
I = \left[ 
\begin{array}{ccc}
\ 1 & 0 & 0 & 0 \\
  0 & 1 & 0 & 0 \\
  0 & 0 & 1 & 0 \\
  0 & 0 & 0 & 1 \\
\end{array}
\right]
\end{equation}

Por el contrario, una matriz inversa es quella matriz \(A^{-1}\), tal que \(AA^{-1} = I\), al multiplicarla por una matriz cualquiera \(A\) da como resultado la matriz identitaria.

\hypertarget{matriz-inversa}{%
\subsection{Matriz inversa}\label{matriz-inversa}}

La utilidad de este concepto se puede ejemplificar con un sistema de una sola ecuación:

\[2x = 10\]

Esta ecuación podemos concebirla como el producto de dos vectores con dimensiones \(1 \times 1\). Para resolverla es necesario multplicar ambos lados de la ecuación por el \textbf{inverso} del coeficiente de \(x\), de modo que:

\[2^{-1} \cdot 2x = 2^{-1} \cdot 10\]

\[x = \frac{1}{5}\]

Por lo tanto, si tenemos que para un sistema lineal la matriz \textbf{cuadrada} \(A\) de coeficientes puede ser invertida como \(A^{-1}\), esta última contiene la solución al sistema mismo:

\[XA = B\]

\[A^{-1}XA = A^{-1}B \rightarrow X = A^{-1}B\]

\hypertarget{definciuxf3n-de-matriz-inversa}{%
\subsubsection{Definción de matriz inversa}\label{definciuxf3n-de-matriz-inversa}}

Una matriz inversa de \(A\) es aquella matriz \(B\) que:

\[AB = BA = I_n\]

al ser multiplicada por \(A\) tiene como resultado la matriz identidad \(I_n\), con dimensiones \(n \times n\). La matriz inversa de \(A\) también recibe, como ya se mencionó el nombre \(A^{-1}\). Cuando se puede obtener la matriz \(A^{-1}\) de \(A\), se dice que esta última es \textbf{invertible} o \textbf{no-singular}, de lo contrario se le llama \textbf{singular}. Cabe resaltar que la matriz inversa \(A^{-1}\) es única para \(A\). Si llegaran a existir dos matrices inversas \(B\) y \(C\) de \(A\), entonces \(B=C\).

\hypertarget{obteniendo-una-matriz-inversa}{%
\subsubsection{Obteniendo una matriz inversa}\label{obteniendo-una-matriz-inversa}}

Para obtener una matriz inversa, también tenemos que resolver un sistema de ecuaciones lineales. Supongamos que \(A\) es una matriz de \(2 \times 2\):

\begin{equation}
A = \left[\begin{array}{ccc}
a_{1, 1} & a_{1, 2} \\
a_{2, 1} & a_{2, 2}
\end{array}
\right]
\end{equation}

Recordemos que \(A\) es la matriz de coeficientes para las variables \(x_1, x_2, \dots, x_n\). Entonces la matriz inversa de \(A\) es:

\begin{equation}
B = \left[\begin{array}{ccc}
b_{1, 1} & b_{1, 2} \\
b_{2, 1} & b_{2, 2}
\end{array}
\right]
\end{equation}

Por lo tanto la ecuación a resolver es:

\begin{equation}
\left[\begin{array}{ccc}
a_{1, 1} & a_{1, 2} \\
a_{2, 1} & a_{2, 2}
\end{array}
\right]
\left[\begin{array}{ccc}
b_{1, 1} & b_{1, 2} \\
b_{2, 1} & b_{2, 2}
\end{array}
\right] =
\left[\begin{array}{ccc}
1 & 0 \\
0 & 1
\end{array}
\right]
\end{equation}

donde las incógnitas son los elementos \(b_{1,1}, \dots b_{n,n}\) de \(B\), la matriz inversa.

\hypertarget{ejemplo-1}{%
\paragraph{Ejemplo}\label{ejemplo-1}}

Tenemos que encontrar la matriz inversa de la matriz de coeficientes:

\begin{equation}
A =
\left[\begin{array}{ccc}
2 & 5 \\
1 & 3
\end{array}
\right]
\end{equation}

Por lo tanto el sistema de ecuaciones a resolver es:

\begin{equation}
\left[\begin{array}{ccc}
2 & 5 \\
1 & 3
\end{array}
\right]
\left[\begin{array}{ccc}
b_{1, 1} & b_{1, 2} \\
b_{2, 1} & b_{2, 2}
\end{array}
\right] =
\left[\begin{array}{ccc}
1 & 0 \\
0 & 1
\end{array}
\right]
\end{equation}

Para ver las ecuaciones a resolver para \(B\), tenemos que multiplicar \(AB\), con lo que obtenemos:

\begin{align}
2 b_{1,1} + 5b_{2,1} = 1\ (F_1) & & 2b_{1, 2} + 5b_{2,2} = 0\ (F_1) \\ \notag
b_{1, 1} + 3b_{2,1} = 0\ (F_2) & & b_{1, 2} + 3b_{2, 2} = 1 (F_2 ) \notag
\end{align}

Utilizando los métodos que ya conocemos, podemos obtener las soluciones:

\begin{equation}
B = \left[ 
\begin{array}
\ b_{1, 1} & b_{1, 2} \\
b_{2, 1} & b_{2, 2}
\end{array}
\right] =
\left[ 
\begin{array}
\ 3 & -5 \\
-1 & 2
\end{array}
\right]
\end{equation}

\hypertarget{ejercicio-3}{%
\paragraph{Ejercicio}\label{ejercicio-3}}

Verifica que las soluciones para \(b_{1,1}, b_{1,2}, b_{2,1}, b_{2,2}\) sean las reportadas y que \(AB = I_2\)

\hypertarget{determinante-de-una-matriz}{%
\subsection{Determinante de una matriz}\label{determinante-de-una-matriz}}

Para continuar con este tema revisaremos rápidamente, qué es y qué hacen los vectores. Por ejemplo, la matriz:

\[
A = \left[ 
\begin{array}{cc}
3 & 0 \\
0 & 2
\end{array}
\right]
\]

puede concebirse como dos vectores \(A_1 = \left[\begin{array}{c}3 \\ 0 \end{array} \right]\) y \(A_2 = \left[\begin{array}{c}0 \\ 2 \end{array} \right]\). Cada uno de estos vectores pueden representarse geométricamente como \emph{indicaciones} de movimiento en el plano cartesiano. El vector \(A_1\) indica \(3\) pasos en dirección \(x\) y \(0\) en dirección \(y\). \(A_2\) indica \(0\) pasos en dirección \(x\) y \(2\) en dirección \(y\).

Cuando representamos estos vectores en el plano vemos que \(A_1\) está alineado con el eje \(x\) y \(A_2\) con el \(y\), y el área formada por ambos corresponde al producto de la primera diagonal de \(A\), \(3 \times 2\) (figura \ref{fig:det}).

\begin{figure}

{\centering \includegraphics[width=41.67in]{Unidad-V/Determinante} 

}

\caption{Representación gráfica del determinante de la matriz A}\label{fig:det}
\end{figure}

Cuando las matrices utilizadas en una transformación lineal (que multiplican a un vector, por ejemplo), no tienen ceros en las casillas fuera de las diagonales, el área proyectada por los vectores será un rombo, por lo que hay que sustraer el área fuera del producto de la diagonal. Por ejemplo la matriz

\[
A = \left[ 
\begin{array}{cc}
1 & 1 \\
\frac{1}{2} & 2
\end{array}
\right]
\]

\begin{figure}

{\centering \includegraphics[width=20.25in]{Unidad-V/Determ-2} 

}

\caption{Representación del área formada por matriz sin ceros fuera de la diagonal}\label{fig:det-2}
\end{figure}

con lo que el área formada por los vectores es \(1 \times 2 - 1/2 \times 1= 3/2\). El área formada por una matriz recibe el nombre de \emph{determinante} (\(\mathrm{det}A\) ó \(|A|\)), puesto que es el efecto neto de la transformación lineal producida por la matriz en el vector a que afecta.

El determinante puede utilizarse con muchos propósitos, uno de ellos es invertir una matriz. Para cualquier matriz \(A\) con dimensiones \(2 \times 2\), \(A\) es invertible si el área que ésta produce es mayor a \(0\).

De manera más general:

\begin{equation}
A = \left[
\begin{array}
\ a_{1,1} & a_{1,2} \\
a_{2,1} & a_{2,2}
\end{array}
\right] \label{eq:mat-A}
\end{equation}

\(A^{-1}\) existe si \(a_{1,1}a_{2,2} - a_{1,2}a_{2,1} \neq 0\) (figura \ref{fig:det-3})

\begin{figure}

{\centering \includegraphics[width=12.81in]{Unidad-V/Deter-3} 

}

\caption{Representación de matriz con determinante de 0.}\label{fig:det-3}
\end{figure}

Para calcular una matriz inversa usando el determinante, necesitamos encontrar la matriz \textbf{adjunta transpuesta}. Para una matriz de \(2 \times 2\) como \eqref{eq:mat-A}:

\begin{equation}
A^{-1} = 
  \frac{1}{a_{1,1}a_{2,2} - a_{1,2}a_{2,1}}
    \left[
      \begin{array}
        \ a_{2,2} & -a_{1,2} \\
        -a_{2,1} & a_{1,1}
      \end{array}
    \right] \label{eq:dos-dos}
\end{equation}

La expresión \(a_{1,1}a_{2,2} - a_{1,2}a_{2,1}\) para una matriz \(A\) de \(2 \times 2\) se conoce como el determinante y se denota como \(\mathrm{det} A\) ó \(|A|\).

Es posible calcular determinantes para cualquier matriz \(A\) con dimensiones \(n \times n\), y en todos los casos, es posible obtener \(A^{-1}\) si \(\mathrm{det}A \neq 0\). Las fórmulas para calcular el determinante de cualquier matriz se van complicando considerablemente conforme aumenta su tamaño, por lo que es necesario utilizar una calculadora de determinantes.

Aquí nos limitaremos a las fórmulas para determinantes de \(2\times 2\) (ecuación \eqref{eq:dos-dos}. La figura \ref{fig:det3x3} muestra el procedimiento para una matrix de \(3 \times 3\).

\begin{figure}

{\centering \includegraphics[width=8.65in]{Unidad-V/determ} 

}

\caption{Esquema del procedimiento para calcular el determinante de una matriz de 3 x 3.}\label{fig:det3x3}
\end{figure}

\hypertarget{valores-propios}{%
\section{Valores propios}\label{valores-propios}}

Para hablar de vectores y valores propios necesitamos definir a un vector como una matriz con una sola columna y \(n\) filas. Todos los ejemplos que veremos aquí son de vectores de \(2 \times 1\).

Los vectores en el sentido en que los trataremos aquí, representan \emph{indicaciones} de movimiento en el plano cartesiano. Por ejemplo el vector:

\[\left[ 
\begin{array}{cc}
2 \\ 3
\end{array}
\right]\]

indica un movimiento de dos posiciones a la derecha en el eje \(x\) y tres posiciones arriba en el eje \(y\) (vector \(\mathbf{u}\) en la figura )

\begin{figure}

{\centering \includegraphics[width=32.06in]{Unidad-V/Vectores} 

}

\caption{Ejemplos de vectores **u** = (2, 3), **u**=(6, 0) y **w** = **u** + **v**.}\label{fig:vectores}
\end{figure}

Las operaciones con vectrores son iguales que las operaciones con matrices, de modo que el vector \(\mathbf{w}\) de la figura \ref{fig:vectores} es:

\[\mathbf{w} =\left[
\begin{array}{c}
2 \\ 3
\end{array}
\right]  + 
\left[
\begin{array}{c}
6 \\ 0
\end{array}
\right] = 
\left[
\begin{array}{c}
2 + 6 \\ 3 +0
\end{array}
\right] =
\left[
\begin{array}{c}
8 \\ 3
\end{array}
\right]\]

Cuando multiplicamos un vector cualquiera por un escalar, éste alargará el vector en relación a su magnitud, tal que :

\[
\mathbf{a} = 2 \times \mathbf{u} = \left[ \begin{array}{c} 4 \\ 6 \end{array} \right]
\]

\begin{figure}

{\centering \includegraphics[width=15.74in]{Unidad-V/Mult-vector} 

}

\caption{Multiplicación de vector por un escalar.}\label{fig:vector-escalar}
\end{figure}

La operación que nos más concierne, sin embargo, en el caso de los valores y vectores propios es la multiplcación de un vector \(\mathbf{x}\) por una matriz \(A\).

Con base en las reglas de multiplicación sabemos que para el vector \(\mathbf{x}\) con dimensiones \(2 \times 1\) y la matriz \(A\) con dimensiones \(2 \times 2\), el producto \(A \mathbf{x}\) tendrá dimensiones de \(2 \times 1\). Dicha multiplicación se conoce como el mapa de \(\mathbf{x}\) y se escribe:

\begin{equation}
\mathbf{x} \mapsto A\mathbf{x} \label{eq:mapa}
\end{equation}

Por ejemplo:

\begin{equation}
    \left[
    \begin{array}{cc}
    3 & 2 \\
    1 & 4
    \end{array}
    \right] \cdot 
        \left[ 
    \begin{array}{c}
    1 \\ 2
    \end{array}
    \right]= 
    \left[
    \begin{array}{c}
    7 \\ 9
    \end{array}
    \right]
\end{equation}

donde \(\mathbf{x}\) es el vector \((1, 2)\), y el mapa \(A \mathbf{x}\) es el vector \((7, 9)\).

El concepto de \textbf{vector y valor propio} está basado en \(\mathbf{x} \mapsto A \mathbf{x}\). Específicamente \(A \mathbf{x}\) es el vector propio y el valor propio es una constante que satisface:

\begin{equation}
A \mathbf{x} = \lambda \mathbf{x} \label{eq:valor-propio}
\end{equation}

Este tipo de operaciones tiene muchas aplicaciones, entre ellas está el resolver sistemas de ecuaciones. Aunque sólo veremos ejemplos con vectores de \(2 \times 1\), todos los ejemplos son generalizables a vectores con cualquier número de filas.

Para calcular \(\lambda\) al igual que hicimos para \(A^{-1}\), comenzamos por resolver la ecuación \eqref{eq:valor-propio}:

\[A \mathbf{x} = \lambda I \mathbf{x} \rightarrow (A - \lambda I) \mathbf{x} = 0\]
La parte \(\lambda I \mathbf{x}\) se debe a que \(I \mathbf{x} = \mathbf{x}\), pero necesitamos una matriz cuadrada para resolver a continuación.

Dado que \(\mathbf{x} \neq 0\), necesitamos que \(A - \lambda I\) sea singular (\(|A| = 0\)), de modo que necesitamos encontrar un valor de \(\lambda\) que satisfaga \(\det (A - \lambda I) = 0\). Entonces si nuestra matriz:

\[A = \left[
\begin{array}{cc}
1 & 2 \\
3 & 2
\end{array}
\right]\]

entonces:

\[A - \lambda I = \left[
\begin{array}{cc}
1 - \lambda & 2 \\
3 & 2 - \lambda
\end{array}
\right]\]

cuyo determinante será:

\[\begin{align}
(1 - \lambda)(2-\lambda) - (2)(3) &= 2 - 3 \lambda +\lambda^2 -6\\
& = \lambda^2 - 3 \lambda -4 \\
& = (\lambda + 1)(\lambda - 4)
\end{align}\]

por lo tanto tenemos las soluciones \(\lambda_1 = -1\) y \(\lambda_2 = 4\). Como es evidente, los vectores propios \(A \mathbf{x}\) que satisfacen los valores propios no son únicos, pero cada valor propio \(\lambda\) está asociado a un vector propio. Para los sistemas lineales con más de dos variables habrá tantos valores propios como dimensiones.

\hypertarget{vectores-propios}{%
\section{Vectores propios}\label{vectores-propios}}

Como vimos anteriormente, los valores propios son independientes de los vectores, es decir, los valores reflejan una propiedad de las matrices. Para encontrar los vectores que satisfacen la propiedad:

\[A \vec{x} = \lambda \vec{x}\]

tenemos que resolver la ecuación para los valores de \(\lambda\). Pongamos por ejemplo la matriz

\[A = \left[\begin{array}{cc} 1 & 2 \\ 3 & 2 \end{array}\right]\]

tiene los valores propios \(\lambda_1 = -1\) y \(\lambda_2 = 4\), y la ecuación a resolver para obtener los vectores propios que corresponden con \(\lambda_1\) son:

\begin{equation}
\left[\begin{array}{cc} 1 & 2 \\ 3 & 2 \end{array}\right] 
\left[\begin{array}{c} x_1 \\ x_2 \end{array} \right] =
-1 \left[\begin{array}{c} x_1 \\ x_2 \end{array} \right]
\end{equation}

pasando el lado izquierdo al derecho, factorizando el vector \(\vec{v}\), y multiplicando \(\lambda\) por la matriz identidad tenemos:

\begin{equation}
\left(
\left[\begin{array}{cc} 1 & 2 \\ 3 & 2 \end{array}\right] +
\left[\begin{array}{cc} 1 & 0 \\ 0 & 1 \end{array} \right] \right) \left[\begin{array}{c} x_1 \\ x_2 \end{array} \right] = \vec{0}
\end{equation}

y sumando las matrices tenemos:

\begin{equation}
\left[\begin{array}{cc} 1 + 1 & 2 \\ 3 & 2 + 1 \end{array}\right] \left[\begin{array}{c} x_1 \\ x_2 \end{array} \right] = \vec{0}
\end{equation}

Una vez encontrado el sistema, podemos resolverlo con el método de la matriz aumentada:

\begin{equation}
\left[ \begin{array}{cc|c}
2 & 2 & 0 \\
3 & 3 & 0
\end{array}\right] 
\begin{array}{c}
(F_1) \\ (F_2)
\end{array}
\end{equation}

aunque si lo resolvemos con el método tradicional \textbf{siempre} llegaremos a la solución \(x_1 = 0; x_2 = 0\). Entonces necesitamos reducir la matrix por fila al modo Echelon. Este método consiste en realizar las operaciones con filas para reducir las filas inferiores a ceros. Por ejemplo, reemplazando \(F_2\) con \(3 F_1 - 2F_2\), y dividiendo \(F_1/2\):

\begin{equation}
\left[ \begin{array}{cc|c}
1 & 1 & 0 \\
0 & 0 & 0
\end{array}\right] 
\begin{array}{c}
(F_3) \\ (F_4)
\end{array}
\end{equation}

Por lo tanto el vector propio para \(\lambda_1 = -1\) es \((1, 1)\). El proceso para encontrar el vector propio para \(\lambda_2\) es el mismo.

\hypertarget{unidad-vi-modelos-lineales-de-poblaciones-estructuradas}{%
\chapter{Unidad VI: Modelos lineales de poblaciones estructuradas}\label{unidad-vi-modelos-lineales-de-poblaciones-estructuradas}}

\hypertarget{proyecciuxf3n-de-matrices-para-modelos-estructurados-la-matriz-de-leslie}{%
\section{Proyección de matrices para modelos estructurados: la matriz de Leslie}\label{proyecciuxf3n-de-matrices-para-modelos-estructurados-la-matriz-de-leslie}}

Los modelos de poblaciones que hemos visto hasta la fecha sólo consideran un estado de desarrollo, tal que toda la población puede ser considerada como \(N\). Al hacer eso,incurrimos en una serie de supuestos biológicos muy fuertes, por ejemplo que todos los individuos nuevos en la población pueden reproducirse en cuanto se incorporan a esta con la misma probabilidad que los individuos más antiguos.

Parte de estos problemas biológicos que el modelo exponencial (\(N_{t+1} = \lambda N_t\)) en tiempo discreto presenta son solucionados por discretización de \(N\), es decir, dividir a \(N\) en una serie de compartimentos que representan estados de desarrollo, de modo que:

\[
N(t) = \left[\begin{array}{c} N_0(t)\\ N_1(t)\\ \vdots \\ N_n(t) \end{array}\right]
\]

donde \(N_0(t)\) representa el cohorte más joven lo neonatos, y \(N_n(t)\) el de mayor edad. Los nacimiento entoces constituyen un flujo de invidiuos nuevos en la población hacia \(N_1(t)\), en función de los cohortes que se pueden reproducir. Para ejemplificar esto, consideremos un organismo cuyas dinámicas ocurren en tiempo discreto, con \(N\) dividido en cuatro edades\(N_0, N_1, N_2\) y \(N_3\). Consideremos primero cómo ocurre la transición entre estos compartimentos, donde sólo sobreviven 40\% de \(N_0\) and \(N_1\), 30\% de \(N_1\) a \(N_2\) y 10\% de \(N_2\) and \(N_3\):

\begin{align}
N_{1, t+1} &= 0.4 N_{0, t} \\
N_{2, t+1} &= 0.3 N_{1, t} \\
N_{3, t+1} &= 0.1 N_{2, t}
\end{align}

\begin{figure}

{\centering \includegraphics[width=8.89in]{Unidad-VI/Modelo-Leslie} 

}

\caption{Diagrama del modelo de proyección de Matrices de Leslie con tres estadíos de desarrollo discretos. Las flechas sólidas representan *flujo* de individuos hacia y desde cada uno de los compartimentos que representa el número de individuos en cada categoría. Las flechas punteadas representan la fuente del flujo de individuos, por ejemplo los Nacimientos en N1 provienen de N2 y N3.}\label{fig:Leslie}
\end{figure}

Para obtener \(N_0(t)\) necesitamos identificar los cohortes o compartimentos que sí se reproducen, entonces si estos son \(N_1\) y \(N_2\), la ecuación para \(N_0\) es:

\begin{align}
N_{0, t+1} &= \mathrm{Crías\ nacidas\ de\ madres} N_{1, t} + \mathrm{Crías\ nacidas\ de\ madres} N_{2, t} \\
N_{0, t+1} &= 2 N_{1, t} + 1.5 N_{2, t}
\end{align}

Claramente, en este modelo estamos asumiendo que la población de interés está constituida únicamente por hembras, de modo que \(2\) y \(1.5\) es el promedio de hembras nacidas por cada hembra en \(N_1\) y \(N_2\). Así tenemos entonces el modelo completo para representar como matriz:

\begin{equation}
\left[\begin{array}{c}
N_{0, t+1} \\ N_{1, t+1} \\ N_{2, t+1} \\ N_{3, t+1}
\end{array}\right] =
\left[ 
\begin{array}{cccc}
0 & 2 & 1.5 & 0 \\
0.4 & 0 & 0 & 0 \\
0  & 0.3& 0 & 0 \\
0  & 0 & 0.1& 0
\end{array}
\right] \label{eq:Leslie}
\end{equation}

La matriz de coeficientes en la ecuación \eqref{eq:Leslie} se conoce omo matriz de Leslie, y la notación corta para dichos sistemas de ecuaciones es:

\[\mathbf{N}(t+1) = L \mathbf{N}(t)\]

\hypertarget{anuxe1lisis-de-los-modelos-de-poblaciones-estructuradas}{%
\section{Análisis de los modelos de poblaciones estructuradas}\label{anuxe1lisis-de-los-modelos-de-poblaciones-estructuradas}}

\hypertarget{tablas-de-vida}{%
\subsection{Tablas de vida}\label{tablas-de-vida}}

Las tablas de vida son una herramienta estadística para organizar datos demográficos de fertilidad y supervivencia de las poblaciones. Existen algunas técnicas para generar tablas de vida, y una de las más confiables son los estudios de cohorte (conjunto de individuos que nacen en un período de tiempo específico, por ejemplo). En estos estudios, se lleva un registro cuidadoso de los sucesos de natalidad y mortalidad (\ref{tab:tabla-vida}).

\begin{table}

\caption{\label{tab:tabla-vida}Tabla de vida de borregos de Dall (*Ovis dalli*) colectada por Olaus Murie en el parque nacional del monte McKinley en Alaska.}
\centering
\begin{tabular}[t]{l|r|r|r}
\hline
Edad & No..Individuos & Muertes & Tasa\\
\hline
0-0.5 & 1000 & 54 & 0.0540\\
\hline
0.5-1 & 946 & 145 & 0.1533\\
\hline
1-2 & 801 & 12 & 0.0150\\
\hline
2-3 & 789 & 13 & 0.0165\\
\hline
3-4 & 776 & 12 & 0.0155\\
\hline
4-5 & 764 & 30 & 0.0393\\
\hline
5-6 & 734 & 46 & 0.0627\\
\hline
6-7 & 688 & 48 & 0.0698\\
\hline
7-8 & 640 & 69 & 0.1078\\
\hline
8-9 & 571 & 132 & 0.2312\\
\hline
9-10 & 439 & 187 & 0.4260\\
\hline
10-11 & 252 & 156 & 0.6190\\
\hline
11-12 & 96 & 90 & 0.9375\\
\hline
12-13 & 6 & 3 & 0.5000\\
\hline
13-14 & 3 & 3 & 1.0000\\
\hline
\end{tabular}
\end{table}

Las tablas de vida entonces, enlistan el número de individuos y el rango de edades al inicio del estudio de cohorte, y el número de muertes que se presentaron mientras los individuos del cohorte se encontraban en cada rango de edades. Este tipo de arreglo permite estimar fácilmente las tasa de mortalidad, como en la tabla \ref{table:tabla-vida}, donde la última columna, es el cociente de la 2a y 3a. Algunas tablas más complejas que la de Olaus Murie contienen una 5a y 6a columnas donde se reporta el número de crías totales y el cociente del número de individuos y crías para cada intervalo de edades.

Las tablas de vida han permitido describir algunas de las diferentes estrategias de los organismos con base en la probabilidad de morir en función de la edad (figura \ref{fig:curva-sup}).

\begin{figure}

{\centering \includegraphics{Mod-mat-ECO_files/figure-latex/curva-sup-1} 

}

\caption{Rleación entre edad y supervivencia en *Ovis dalli*, de acuerdo con la tabla de vida del ejemplo.}\label{fig:curva-sup}
\end{figure}

La diversidad de estrategias reproductivas existentes resultan en curvas de supervivencia con otras formas. Por ejemplo la curva de \emph{Ovis dalli} es convexa, lo cual quiere decir que la tasa de mortalidad aumenta con la edad. En plantas sin embargo, la curva es cóncava y quiere decir que las muertes ocurren mayoritariamente a edades tempranas, mientras que una pequeña fracción de individuos sobrevive a edades muy avanzadas.

Estras estrategias reproductivas también se manifiestan en la estructura de edades dentro de las poblaciones, es decir la cantidad de individuos neonatos, jóvenes y adultos que hay.

\hypertarget{la-ecuaciuxf3n-de-euler}{%
\subsection{La ecuación de Euler}\label{la-ecuaciuxf3n-de-euler}}

En demografía, se le conoce como la ecuación característica, y se utiliza para calcular la tasa neta reproductiva de una población. La interpretación del número producido por la ecuación a partir de los datos contenidos en una tabla de vida es \textbf{el número de hembras con el que cada hembra es reemplazada}, y este número se conoce como \(R_0\). Para calularlo necesitamos una tabla de vida para un cohorte de hembras, y hacer la suma:

\begin{equation}
R_0 = \sum l_x m_x \label{eq:Euler-demog}
\end{equation}

donde \(l_x\) es la proporción viva del cohorte en cada categoría de edades (1 - columna 4 de tabla \ref{tab:tabla-vida}, o tasa de supervivencia por edad), y \(m_x\), el número promedio de descendientes hembras por cada hembra en cada categoría de edades \(x\).

Para dar un ejemplo utilizaremos otra tabla de vida para ardillas terrestres de Belding \citep{Sherman1984}:

\begin{table}

\caption{\label{tab:tabla-belding}Tabla de vida para hembras de *Spermophilus beldingi*.}
\centering
\begin{tabular}[t]{r|r|r}
\hline
Edad & l\_x & m\_x\\
\hline
0 & 1.000 & 0.000\\
\hline
1 & 0.386 & 1.040\\
\hline
2 & 0.197 & 2.171\\
\hline
3 & 0.106 & 2.450\\
\hline
4 & 0.054 & 3.113\\
\hline
5 & 0.029 & 1.875\\
\hline
6 & 0.014 & 1.575\\
\hline
7 & 0.008 & 2.000\\
\hline
8 & 0.006 & 1.440\\
\hline
9 & 0.002 & 1.632\\
\hline
\end{tabular}
\end{table}

Como es evidente, nos queda por calcular \(l_x m_x\) que corresponde al producto de las 2a y 3a columnas. Una vez obtenido el producto podemos hacer la sumatoria, el código de \textbf{R} para hacer estas operaciones es:

\begin{Shaded}
\begin{Highlighting}[]
\NormalTok{lxmx }\OtherTok{\textless{}{-}}\NormalTok{ tab2}\SpecialCharTok{$}\NormalTok{l\_x}\SpecialCharTok{*}\NormalTok{tab2}\SpecialCharTok{$}\NormalTok{m\_x}
\FunctionTok{sum}\NormalTok{(lxmx)}
\end{Highlighting}
\end{Shaded}

\begin{verbatim}
## [1] 1.361258
\end{verbatim}

lo que indica que cada hembra en promedio será reemplazada por 1.36 hembras.

El concepto de \(R_0\) es también muy utilizado en epidemiología, donde tiene las mismas implicaciones. Si cada hembra es reemplazada en promedio por otra hembra (\(R_0 \approx 1\)), la población permanecerá estable, por el contrario si \(R_0 < 1\) la población tenderá a extinguirse. En epidemiología cada infectado deberá contagiar a más de uno para que haya una epidemia.

\hypertarget{valor-reproductivo-y-autovectores}{%
\subsection{Valor reproductivo y autovectores}\label{valor-reproductivo-y-autovectores}}

Las matrices de coeficientes de transición entre estados de desarrollo y las tasas de natalidad asociadas a cada estado pueden ser representadas con menos parámetros puesto que el crecimiento de la población continúa siendo exponencial. La manera de representar toda la matriz de coeficientes como uno sólo es con los valores propios. Veamos el ejemplo de la ecuación \eqref{eq:Leslie}, cuyos valores propios son:

\[\lambda = \{0, 0.99, -0.24, -0.75 \}\]

Algunos de estos valores propios, como los negativos, son evidentemente, biológicamente insignificantes. En la práctica se elije sólamente uno de los valores propios, denominado \textbf{valor dominante} y que corresponde al valor más grande.

Como ya vimos en la unidad de álgebra de matrices, los valores propios también pueden ser números complejos, pero las implicaciones biológicas de estos valores propios no forman parte de esta unidad de modelos lineales.

\hypertarget{aplicaciones-de-lo-modelo-de-matrices}{%
\section{Aplicaciones de lo modelo de matrices}\label{aplicaciones-de-lo-modelo-de-matrices}}

Los modelos con matriz de Leslie tienen muchas aplicaciones en ecología, debido a que son sumamente flexibles y relativamente fáciles de implementar. Para comprender primero cómo funcionan y el tipode dinámicas que surgen con estos modelos echemos un vistazo al \href{Unidad-VI/Modelos-matriz.R}{script de R}

  \bibliography{book.bib,packages.bib,citas.bib}

\end{document}
